\documentclass[12pt,a4paper]{report}

% ============================================================
% PAGE LAYOUT
% ============================================================

\usepackage[
  margin=1in
]{geometry}

% ============================================================
% CORE GRAPHICS, TABLES AND MATHEMATICS
% ============================================================

\usepackage{graphicx}
\usepackage{booktabs}
\usepackage{array}
\usepackage{longtable}
\usepackage{tabularx}
\usepackage{colortbl}
\usepackage{multirow}

\usepackage{amsmath}
\usepackage{amssymb}

\usepackage{float}

% ============================================================
% FIGURE CAPTIONS AND SUBFIGURES
% ============================================================

\usepackage{caption}
\usepackage{subcaption}

% ============================================================
% DIAGRAMS
% ============================================================

% All diagrams used by the report are pre-rendered PDF or PNG files.
% TikZ and PGFPlots are intentionally not loaded here because loading
% those packages substantially increases Overleaf start-up time.

% ============================================================
% REFERENCES
% ============================================================

\usepackage[
  round,
  authoryear
]{natbib}

% ============================================================
% TEXT FORMATTING AND COLOUR SUPPORT
% ============================================================

\usepackage{setspace}
\usepackage{enumitem}
\usepackage{fancyhdr}
\usepackage{xcolor}

\onehalfspacing

% Provides additional flexibility when LaTeX is resolving
% difficult line breaks.
\emergencystretch=2em

% ============================================================
% PROJECT COLOURS
% ============================================================

% Main project colour, adapted from the viva presentation.
\definecolor{projectblue}{RGB}{28,76,128}

% Light alternating table row.
\definecolor{tableStripe}{RGB}{239,245,250}

% Highlight for selected or winning model.
\definecolor{winnerRow}{RGB}{218,237,224}

% Pipeline stage colours.
\definecolor{stageone}{RGB}{224,237,249}
\definecolor{stagetwo}{RGB}{224,241,228}
\definecolor{stagethree}{RGB}{250,239,214}

% Additional pale background colours.
\definecolor{lightblue}{RGB}{232,240,248}
\definecolor{verylightblue}{RGB}{246,249,252}


% ============================================================
% BODY TYPOGRAPHY
% ============================================================

% The Singapore-blue visual treatment is intentionally limited to the
% title page. The dissertation body uses the standard report-class
% chapter and section styling for a conventional academic appearance
% and faster compilation.

\captionsetup{
  font=small,
  labelfont=bf,
  justification=centering,
  singlelinecheck=false
}

% ============================================================
% CLICKABLE LINKS, CONTENTS AND PDF BOOKMARKS
% ============================================================

% Hyperref is loaded only after projectblue has been defined.
% Every link is explicitly blue: no red text and no red boxes.
\usepackage{xurl}
\usepackage[
  pdfencoding=auto,
  psdextra
]{hyperref}
\usepackage{bookmark}

\hypersetup{
  colorlinks=true,
  linkcolor=projectblue,
  citecolor=projectblue,
  urlcolor=projectblue,
  filecolor=projectblue,
  menucolor=projectblue,
  anchorcolor=projectblue,
  linktoc=all,
  pdfborder={0 0 0},
  bookmarks=true,
  bookmarksopen=true,
  bookmarksnumbered=true,
  bookmarksopenlevel=2,
  pdfpagemode=UseOutlines,
  pdftitle={
    Restricted Dog Breed Classification Using Deep Learning
    and Public Image Datasets
  },
  pdfauthor={Ronan Downes},
  pdfsubject={
    CNN transfer learning for restricted and unrestricted
    dog-image classification
  },
  pdfkeywords={
    convolutional neural networks,
    transfer learning,
    dog-breed classification,
    InceptionResNetV2,
    Grad-CAM,
    deep learning
  }
}


% ============================================================
% CLICKABLE PAGE NUMBERS: RETURN TO CONTENTS
% ============================================================

% Clicking a footer page number returns to the actual Contents page (Roman page iv).
% Headers and page-number typography remain plain and black.
\pagestyle{fancy}
\fancyhf{}
\fancyhead[L]{\nouppercase{\leftmark}}
\fancyfoot[C]{\hypersetup{linkcolor=black}\hyperlink{page.iv}{\thepage}\hypersetup{linkcolor=projectblue}}
\renewcommand{\headrulewidth}{0.4pt}
\renewcommand{\footrulewidth}{0pt}

\fancypagestyle{plain}{%
  \fancyhf{}%
  \fancyfoot[C]{\hypersetup{linkcolor=black}\hyperlink{page.iv}{\thepage}\hypersetup{linkcolor=projectblue}}%
  \renewcommand{\headrulewidth}{0pt}%
  \renewcommand{\footrulewidth}{0pt}%
}

% ============================================================
% TABLE SETTINGS
% ============================================================

% Increases vertical spacing in tables.
\renewcommand{\arraystretch}{1.35}

% Default horizontal space between table columns.
\setlength{\tabcolsep}{6pt}

% Reusable column types.
%
% L = left-aligned flexible-width column
% C = centred flexible-width column
% R = right-aligned flexible-width column
%
\newcolumntype{L}{>{\raggedright\arraybackslash}X}
\newcolumntype{C}{>{\centering\arraybackslash}X}
\newcolumntype{R}{>{\raggedleft\arraybackslash}X}

% Fixed-width left, centre and right columns.
\newcolumntype{P}[1]{
  >{\raggedright\arraybackslash}p{#1}
}

\newcolumntype{Q}[1]{
  >{\centering\arraybackslash}p{#1}
}

\newcolumntype{T}[1]{
  >{\raggedleft\arraybackslash}p{#1}
}

% ============================================================
% FIGURE AND TABLE CAPTION SETTINGS
% ============================================================


\captionsetup[table]{
  position=above,
  skip=8pt
}

\captionsetup[figure]{
  position=below,
  skip=6pt
}

% ============================================================
% SECTION AND CONTENTS DEPTH
% ============================================================

% Number subsections down to the third level.
\setcounter{secnumdepth}{3}

% Include subsections down to the third level in the contents.
\setcounter{tocdepth}{2}

% ============================================================
% GRAPHICS DIRECTORY
% ============================================================

\graphicspath{
  {figures/}
}

% ============================================================
% REUSABLE RESULT COMMANDS
% ============================================================

% Highlight the selected model name.
\newcommand{\selectedmodel}[1]{
  \textbf{#1}
}

% Typeset percentages consistently.
\newcommand{\resultpercent}[1]{#1\%}

% Typeset a reduction or increase in percentage points.
\newcommand{\percentagepoints}[1]{
  #1 percentage points
}

% ============================================================
% CONDITIONAL BUILD SETTINGS
% ============================================================

% This file should define the draft-image switch, for example:
%
% \newif\ifdraftimages
% \draftimagesfalse
%
% OVERLEAF FREE-PLAN BUILD SWITCH
% --------------------------------
% Normal final build: figures are rendered.
\newif\ifdraftimages
\draftimagesfalse

% During text-only editing, change the line above to:
% \draftimagestrue
% This replaces figures with labelled boxes and compiles faster.


% When draft-image mode is enabled, images are represented
% by empty bounding boxes. This greatly reduces compilation time.
\ifdraftimages
  \setkeys{Gin}{draft}
\fi

% ============================================================
% DOCUMENT INFORMATION
% ============================================================

\title{
  Restricted Dog Breed Classification Using Deep Learning
  and Public Image Datasets
}

\author{
  Ronan Downes\\[1cm]
  Supervisor: Sam Weiss\\
  Course Director: David McQuaid\\[1cm]
  MSc in Data Analytics\\
  CCT College Dublin
}

\date{July 2026}

% ============================================================
% DOCUMENT
% ============================================================

\begin{document}

% ============================================================
% TITLE PAGE
% ============================================================

\begin{titlepage}
  \thispagestyle{empty}
  \centering

  \vspace*{1.2cm}
  {\color{projectblue}\rule{\textwidth}{2.2pt}}\\[1.3cm]

  {\sffamily\bfseries\color{projectblue}\fontsize{29}{34}\selectfont
   Restricted Dog Breed Classification\\[0.25em]
   Using Deep Learning and\\[0.25em]
   Public Image Datasets\par}

  \vspace{0.8cm}
  {\color{lightblue}\rule{0.72\textwidth}{1.1pt}}\\[1.15cm]

  {\Large\sffamily Ronan Downes\par}

  \vfill

\begin{minipage}{0.72\textwidth}
  \centering
  {\large\sffamily
    Supervisor: Sam Weiss\\[0.35em]
    Course Director: David McQuaid\\[1.1em]
    MSc in Data Analytics\\
    CCT College Dublin\\[1.1em]
    \textbf{Word count: 14,030}
  }
\end{minipage}
  \vfill

  {\large\sffamily\color{projectblue}July 2026\par}
  \vspace{0.8cm}
  {\color{projectblue}\rule{\textwidth}{2.2pt}}
\end{titlepage}
\pdfbookmark[0]{Title Page}{titlepage}

% ============================================================
% FRONT MATTER
% ============================================================

\pagenumbering{roman}

% ------------------------------------------------------------
% DECLARATION
% ------------------------------------------------------------

\chapter*{Declaration}
\addcontentsline{toc}{chapter}{Declaration}

I declare that this dissertation is my own work, except where the
work of others is acknowledged through citation and reference.
The experimental notebooks, saved outputs, figures and tables
supplied with the project form the reproducibility record for
the study.

% ------------------------------------------------------------
% ACKNOWLEDGEMENTS
% ------------------------------------------------------------

\chapter*{Acknowledgements}
\addcontentsline{toc}{chapter}{Acknowledgements}

I would like to thank my supervisor, Sam Weiss, for his guidance
and support throughout the capstone project. I also thank David
McQuaid, Course Director of the MSc in Data Analytics, for his
direction of the programme and his guidance on the structure,
evaluation and presentation of the thesis.

I am grateful to the teaching staff of CCT College Dublin for
the technical and academic foundation provided across the
programme. Finally, I thank my wife Airida and my children
Gabija and Erik for their patience, encouragement and
understanding throughout the work.

% ------------------------------------------------------------
% ABSTRACT
% ------------------------------------------------------------

\chapter*{Abstract}
\addcontentsline{toc}{chapter}{Abstract}

This dissertation investigates whether transfer learning with
convolutional neural networks can support the visual
classification of dog breeds subject to additional controls
under Irish law.

A binary restricted/unrestricted image-classification task was
constructed from the Stanford Dogs dataset using six represented
restricted breeds and a reproducibly selected sample of
twenty-four unrestricted breeds. Six ImageNet-pretrained
architectures were evaluated as frozen feature extractors under
a common pipeline. The models compared were VGG16, ResNet50,
InceptionV3, Xception, InceptionResNetV2 and NASNetMobile.

InceptionResNetV2 provided the strongest frozen result, achieving
97.08\% accuracy, 87.38\% precision, 98.42\% recall and an F1
score of 92.57\%. Fine-tuning the final twenty layers improved
accuracy to 98.25\%, precision to 93.88\% and F1 to 95.34\%,
while recall reduced slightly to 96.84\%. The corresponding
confusion matrix contained 826 true negatives, 12 false
positives, 6 false negatives and 184 true positives.

The findings therefore demonstrate that fine-tuning changed the
operating trade-off rather than uniformly improving every
metric: false positives were more than halved, but false
negatives increased from three to six.

A qualitative Grad-CAM study of twelve fixed cases was used to
examine true-positive, true-negative, false-positive and
false-negative decisions. The maps frequently overlapped the
dog, but some predictions were influenced by pose, scale,
contrast and background structure.

The study concludes that transfer learning is effective for this
controlled dataset, while emphasising that a photograph-only
classifier cannot establish legal breed identity and should not
be treated as an autonomous enforcement mechanism. The main
contribution is a reproducible comparison of established CNN
architectures, a metric-aware fine-tuning analysis, and a
critical account of the technical, legal and ethical limits of
image-based restricted-breed classification.


% CONTENTS, FIGURES AND TABLES
% ============================================================

\phantomsection
\hypertarget{contentspage}{}
\tableofcontents
\listoffigures
\listoftables

\clearpage

% ============================================================
% MAIN MATTER
% ============================================================

\pagenumbering{arabic}

\chapter{Introduction}
\label{chap:introduction}

\section{Background and motivation}
Dog-breed recognition appears at first to be an ordinary image-classification problem. A photograph is supplied, a model extracts visual information, and a breed label is returned. The practical and ethical context of this dissertation makes the task considerably more complex. Irish dog-control policy applies additional controls to named breeds and to strains or crosses of those breeds. The controlled list includes the American Pit Bull Terrier, English Bull Terrier, Staffordshire Bull Terrier, Bull Mastiff, Doberman Pinscher, German Shepherd, Rhodesian Ridgeback, Rottweiler, Japanese Akita, Japanese Tosa and dogs commonly described as Bandogs. In public, these dogs are subject to requirements that include secure muzzling and handling on a strong lead no longer than two metres by a person over sixteen who is capable of controlling the animal. Separate regulations introduced for XL Bully type dogs prohibit importation, breeding, sale and rehoming, and require a Certificate of Exemption for continued ownership of existing animals. The legal category is therefore consequential: an identification can influence the obligations imposed on an owner and the treatment of the dog.

The law does not, however, turn a visually heterogeneous group into a coherent biological image class. A German Shepherd, a Rhodesian Ridgeback and a Staffordshire Bull Terrier differ substantially in coat, silhouette, head shape and body proportions. The positive class in this project is united by regulation, not by a single shared morphology. Conversely, unrestricted breeds can resemble restricted breeds in particular views. A Black-and-Tan Coonhound may share colour cues with a Rottweiler; a Kelpie may resemble a lean shepherd-type dog; and several mastiff or bull-type breeds can be difficult to distinguish from one photograph. The task therefore combines fine-grained recognition with a heterogeneous target class and asymmetric practical consequences.

Fine-grained visual categorisation requires a model to distinguish subcategories within the same broad object class. The image almost always contains a dog; the problem is not to separate dogs from cars or buildings, but to identify comparatively subtle combinations of muzzle length, ear position, coat texture, body mass, stance, colouring and facial structure. The Stanford Dogs dataset was designed as a benchmark for this type of classification and contains images from 120 breeds \citep{khosla2011novel,stanforddogs}. It provides a useful public source for a controlled academic experiment, but it was not designed for legal identification. It contains curated internet images, uneven backgrounds, variable viewpoints and breed labels inherited from source data. It does not contain every legally relevant breed or type, and it cannot represent the full diversity of crossbred dogs encountered in practice.

The central motivation is therefore not to automate a legal decision. It is to investigate what modern transfer learning can and cannot achieve under clearly stated conditions. Computer vision models can provide repeatable numerical outputs, but repeatability is not equivalent to legal validity. A model can learn correlations in a dataset without learning a stable concept of breed, and a highly accurate test result can coexist with brittleness under distribution shift. This dissertation treats the classifier as an experimental artefact and evaluates it critically as such.

\section{Research problem}
The empirical problem is binary classification. Given an image drawn from the prepared test set, the model must estimate the probability that it belongs to the restricted class. The positive class consists of six restricted breeds represented in Stanford Dogs: Bull Mastiff, Doberman Pinscher, German Shepherd, Rhodesian Ridgeback, Rottweiler and Staffordshire Bull Terrier. The negative class consists of twenty-four unrestricted breeds selected using a fixed random seed. The resulting task asks whether a shared classifier can learn a boundary around a legally defined but visually diverse positive group.

Three technical difficulties follow. First, the class distribution is imbalanced: unrestricted images outnumber restricted images. A model that predicts the majority class too readily could report an attractive accuracy while missing restricted examples. Second, the positive class is multimodal. The network must accept several distinct visual patterns as restricted, rather than learning one breed prototype. Third, training data are limited relative to the capacity of modern CNNs. Training a large architecture from random initialisation would create a substantial risk of overfitting and would require more computational resources than were available. These constraints motivate transfer learning, metric-aware evaluation and careful checkpointing.

The policy context adds a fourth difficulty: the costs of errors are not interchangeable. A false negative labels a restricted image as unrestricted; a false positive labels an unrestricted image as restricted. In a hypothetical screening setting, false negatives could fail to flag a relevant case, while false positives could subject an owner and animal to unwarranted scrutiny. This is why the dissertation reports precision, recall, F1, area under the ROC curve, area under the precision--recall curve, Matthews correlation coefficient and confusion-matrix counts rather than relying on accuracy alone.

\section{Research aim, questions and objectives}
The overall aim is to benchmark and critically evaluate a reproducible ImageNet transfer-learning pipeline for binary restricted-versus-unrestricted dog-image classification using public data. The benchmark compares six pretrained convolutional neural networks under one controlled experimental design before limited fine-tuning of the strongest frozen model.

The research questions are:
\begin{enumerate}
    \item How well do six frozen ImageNet-pretrained CNNs distinguish the constructed restricted and unrestricted classes under one common pipeline?
    \item Which architecture provides the strongest balance of precision, recall, F1, PR-AUC and MCC?
    \item What changes after limited fine-tuning of the strongest frozen model, particularly in the false-positive/false-negative trade-off?
    \item What do Grad-CAM visualisations reveal about the image regions associated with correct and incorrect predictions?
    \item Why can the experimental classifier not be treated as an autonomous legal breed-identification system?
\end{enumerate}

The operational objectives are to construct a reproducible binary dataset; benchmark VGG16, ResNet50, InceptionV3, Xception, InceptionResNetV2 and NASNetMobile; select a model using metrics suited to imbalance; fine-tune the selected architecture; preserve logs, checkpoints and predictions; analyse representative errors; and relate the findings to the limitations of breed-specific regulation and high-stakes automated classification.

\section{Scope and boundaries}
The project is deliberately narrower than real-world breed identification. It classifies images already assumed to show one dog from one of the selected Stanford categories. It does not detect dogs within unrestricted scenes, identify mixed ancestry, estimate behaviour or dangerousness, verify legal documents, or make decisions about seizure or ownership. The label ``restricted'' is a project label derived from a legal list, not an assertion that an individual animal is dangerous.

The study also distinguishes legal status from visual classifiability. Some legally named breeds are absent from Stanford Dogs, including the Japanese Tosa and the non-standardised Bandog type. The American Pit Bull Terrier is not represented as a straightforward Stanford category, and the XL Bully is outside the dataset's original taxonomy. The model cannot reliably identify categories for which it has no training representation. This is not a minor implementation gap; it demonstrates a general principle of supervised learning: a model's competence is bounded by its data, labels and deployment distribution.

\section{Contribution}
The dissertation contributes a complete comparative workflow rather than a claim of production readiness. It provides a controlled benchmark of VGG16, ResNet50, InceptionV3, Xception, InceptionResNetV2 and NASNetMobile; a transparent rule for choosing the strongest frozen model; a fine-tuning comparison using the same held-out test set; and a qualitative explanation study built around fixed true-positive, true-negative, false-positive and false-negative cases. The notebooks save intermediate artefacts so that results can be inspected rather than accepted solely from narrative description.

A second contribution is critical interpretation. The final model achieved high performance, but the report does not convert this result into a claim that image-based legal identification is solved. Instead, it examines how class definition, dataset composition, threshold choice, background cues, explanation limitations and asymmetric error costs constrain interpretation. This approach aligns with the Level 9 requirement to integrate technical knowledge with methodological justification, reflective evaluation and ethical analysis.

\section{Dissertation structure}
Chapter~\ref{chap:cnn} explains the foundations of CNN image classification, including convolution, activation, pooling, normalisation, transfer learning and fine-tuning. Chapter~\ref{chap:literature} reviews the architecture families and related literature. Chapter~\ref{chap:methodology} documents the experimental design and reproducibility controls. Chapter~\ref{chap:results} reports the frozen model-zoo and fine-tuning results. Chapter~\ref{chap:explanation} analyses the Grad-CAM cases. Chapter~\ref{chap:discussion} provides critical evaluation, ethics, limitations and future work. Chapter~\ref{chap:conclusion} concludes the study.

\section{Why this is a data-analytics problem}
The project sits within data analytics rather than pure software engineering because the principal challenge is not merely to make code run. It is to define a valid target, prepare and audit data, select measurements, compare alternatives and interpret uncertainty. The label engineering step translates a policy category into a machine-learning target. That translation is itself an analytical decision. A technically correct training loop would still produce weak research if the class definition, split integrity or evaluation metrics were poorly justified.

The project also illustrates the distinction between prediction and explanation. Predictive modelling asks whether patterns in the input support accurate held-out classification. Explanatory or causal analysis asks why a phenomenon occurs and whether changing one factor would change the outcome. This dissertation performs prediction. It does not establish that visual breed features cause dangerous behaviour or that the legal list is an optimal public-safety intervention. Keeping these questions separate prevents a successful classifier from being used as evidence for a claim it was never designed to test.

A further analytical issue is aggregation. The restricted class combines several breeds, while reported metrics combine all test images. Aggregate performance can hide local weakness. A model may be excellent for a visually distinctive breed and poor for another, yet still achieve a high overall score. This motivates breed-level analysis in future work and explains why the Grad-CAM chapter examines individual cases rather than relying solely on summary statistics.

\section{Terminology}
Throughout the dissertation, ``restricted'' refers to the positive project label derived from the Irish controlled-breed list and the six relevant Stanford categories. ``Unrestricted'' refers to the selected negative categories; it does not mean that the dog is exempt from all general dog-control law. ``Breed classification'' refers to image-based statistical prediction, not genetic verification. ``Fine-tuning'' means updating a limited part of a pretrained backbone after the frozen-head stage. ``Explanation'' refers to a post-hoc visual diagnostic and not to a legally sufficient account of a decision.


\chapter{Convolutional Neural Network Foundations}
\label{chap:cnn}

\section{From pixels to learned representations}
A digital colour image can be represented as a three-dimensional tensor. The first two dimensions describe spatial position and the third describes colour channels, commonly red, green and blue. An input of $299\times299\times3$ therefore contains 268,203 numerical values before batching. A conventional fully connected network would connect every input value to every neuron in the next layer. Even a modest dense layer would create an excessive number of parameters and would ignore the fact that neighbouring pixels are spatially related. CNNs address these problems through local connectivity and weight sharing \citep{lecun1998gradient,lecun2015deep}.

A learned image representation is hierarchical. Early layers often respond to local contrasts, oriented edges and colour transitions. Intermediate layers combine those responses into textures, corners and parts. Later layers can represent higher-level configurations that support the classification objective. These descriptions are helpful intuitions rather than fixed rules: the actual representation depends on architecture, optimisation, data and loss. Nonetheless, the hierarchy explains why features learned from a large generic dataset can be useful for a smaller specialised task.

\section{Convolution}
A convolutional layer applies a bank of learned kernels to local regions of its input. For an input feature map $X$ and kernel $K$, a simplified two-dimensional cross-correlation operation can be written as
\begin{equation}
Y(i,j)=b+\sum_{u=0}^{h-1}\sum_{v=0}^{w-1}K(u,v)X(i+u,j+v),
\end{equation}
where $h$ and $w$ are the kernel dimensions and $b$ is a learned bias. Deep-learning libraries commonly implement cross-correlation while retaining the conventional term ``convolution''. For colour images and intermediate tensors, the sum also extends across input channels.

The kernel slides across the image using the same weights at each location. This weight sharing is fundamental. A vertical-edge detector learned in the upper-left region can respond to the same pattern elsewhere. The layer produces a feature map indicating where the learned pattern is strongly expressed. Multiple filters produce multiple maps, allowing a layer to detect different visual properties.

\begin{figure}[H]
\centering
\includegraphics[width=0.85\textwidth]{convolution_diagram.pdf}
\caption{Conceptual convolution operation. A learned kernel is applied locally and produces one activation at each position.}
\label{fig:convolutionconcept}
\end{figure}

Kernel size determines the immediate receptive field of a unit. A $3\times3$ kernel observes a neighbourhood of nine spatial positions, while deeper stacks of such layers acquire progressively larger effective receptive fields. This is one reason modern architectures can model global shape while using predominantly small kernels. The values of the kernels are not hand-selected edge templates in the final model. They are trainable parameters adjusted by gradient-based optimisation to reduce the classification loss.

\section{Stride, padding and receptive field}
Stride describes how far the kernel moves between applications. A stride of one evaluates adjacent locations; a stride of two reduces spatial resolution because every second location is sampled. Padding controls treatment of image borders. ``Valid'' convolution applies a kernel only where it fully overlaps the input, reducing width and height. ``Same'' padding adds border values, usually zeros, so that stride-one convolution preserves spatial dimensions.

The receptive field of a deeper unit is the portion of the original input that can influence it. Receptive field grows with depth, kernel size, stride and pooling. In breed classification, a small receptive field may encode coat texture or an ear boundary, whereas a large receptive field can combine head, torso and stance. The distinction matters because breed cues occur at multiple scales. Inception architectures explicitly process several scales in parallel, while residual networks construct very deep hierarchies that enlarge contextual coverage.

\section{Non-linearity and ReLU}
A sequence of purely linear operations collapses to a single linear transformation, regardless of depth. Neural networks therefore apply nonlinear activation functions. The rectified linear unit is
\begin{equation}
f(x)=\max(0,x).
\end{equation}
ReLU retains positive activations and sets negative activations to zero. Its simplicity and favourable gradient behaviour contributed to the successful training of deep image networks \citep{nair2010rectified,krizhevsky2012imagenet}. It also creates sparse activation patterns in which a feature is active only for particular inputs.

ReLU should not be interpreted as proving that a visual feature is present. It is one element in a distributed representation. Later layers combine many activations, and the same unit may participate in different decisions. This distributed character is one reason that post-hoc explanation methods such as Grad-CAM provide coarse localisation rather than a complete causal account.

\section{Pooling and spatial reduction}
Pooling summarises a local neighbourhood of a feature map. In $2\times2$ max pooling, the largest value in each local window is retained. For example,
\begin{equation}
\begin{bmatrix}1&3\\2&7\end{bmatrix}\longrightarrow 7.
\end{equation}
The operation reduces spatial dimensions and preserves the strongest local response. It can provide limited tolerance to small translations: if a strong activation moves slightly within the window, the pooled value may remain unchanged. Average pooling instead retains the mean response.

Pooling is not synonymous with learning. A max-pooling layer has no trainable kernel values; it applies a fixed reduction. Modern architectures sometimes replace explicit pooling with strided convolution or specialised reduction blocks. InceptionV3 and InceptionResNetV2 use mixed reduction modules, while Xception uses strided separable convolutions in parts of its flow. The architectural purpose is similar: reduce resolution while increasing semantic abstraction and channel depth.

At the end of many contemporary CNNs, global average pooling computes one mean value per feature channel. If a final tensor has shape $8\times8\times1536$, global average pooling produces a vector of length 1536. This avoids the extremely large parameter count created by flattening every spatial position into a dense layer. Global average pooling was proposed as a structural alternative to conventional dense classification heads and can improve correspondence between feature channels and categories \citep{lin2013network}.

\section{Fully connected layers and binary output}
A dense or fully connected layer forms weighted combinations of all values in its input vector. In a binary classifier, the final layer commonly contains one unit with sigmoid activation:
\begin{equation}
\sigma(z)=\frac{1}{1+e^{-z}}.
\end{equation}
The output lies between zero and one and is interpreted as a model score for the positive class. It is not automatically a calibrated probability. A decision threshold, commonly 0.5, converts the continuous score into a class prediction. Moving the threshold changes the balance of false positives and false negatives.

The binary cross-entropy loss for target $y\in\{0,1\}$ and score $p$ is
\begin{equation}
\mathcal{L}(y,p)=-\left[y\log(p)+(1-y)\log(1-p)\right].
\end{equation}
Confident incorrect predictions incur a large penalty. During training, gradients of the loss are propagated backward through the classification head and any trainable backbone layers. The optimiser updates weights in a direction expected to reduce future loss.

\section{Backpropagation and optimisation}
Backpropagation applies the chain rule to compute the contribution of each trainable parameter to the loss. It is not a separate learning algorithm from gradient descent; it is an efficient method for obtaining gradients in a composed computational graph. The optimiser then decides how to use these gradients. Adam maintains adaptive estimates of first and second moments of the gradient and is widely used because it provides effective optimisation across parameters with different gradient scales \citep{kingma2014adam}.

Learning rate controls update magnitude. A rate that is too large can destabilise training or destroy useful pretrained features. A rate that is too small can make progress extremely slow. Fine-tuning therefore normally uses a much lower learning rate than training a newly initialised classification head. In this project, the frozen stage learned only the new head, whereas fine-tuning updated a limited number of upper backbone layers at $10^{-5}$.

\section{Batch normalisation}
Batch normalisation standardises intermediate activations using mini-batch statistics and then applies learned scale and shift parameters. It was introduced to improve optimisation stability and permit more aggressive learning rates \citep{ioffe2015batch}. During inference, moving estimates accumulated during training are used rather than the current batch's statistics.

Pretrained architectures contain many batch-normalisation layers. Their handling during fine-tuning requires care. Updating batch statistics with small or unrepresentative batches can disturb a pretrained representation. For this reason, transfer-learning implementations often keep the base model in inference mode while selectively enabling trainable weights in upper layers. The exact behaviour depends on the framework and model construction, and it is more precise to describe the implemented code than to make a general claim that every unfrozen batch-normalisation layer behaves identically.

\section{Dropout and regularisation}
Dropout randomly sets a proportion of activations to zero during training and rescales the remaining activations so that expected magnitude is maintained. The method discourages units from depending excessively on particular co-adapted features \citep{srivastava2014dropout}. During inference, dropout is disabled. A dropout rate is therefore not the proportion of model parameters permanently removed; it is a stochastic training mechanism.

Regularisation also arises from data augmentation, early stopping, pretrained initialisation, limited fine-tuning and model selection based on validation data. These mechanisms address different aspects of overfitting. Dropout regularises the classification head, augmentation alters the observed input distribution, and early stopping limits the number of updates after validation performance ceases to improve.

\section{Data augmentation}
The project used light image augmentation in the training graph, including horizontal flipping, small rotations and zoom transformations. Augmentation exposes the classifier to plausible variations without changing the semantic label. A dog remains the same breed when it is mirrored horizontally or framed slightly differently. The objective is not to create new biological evidence, but to reduce reliance on accidental image coordinates and exact framing.

Augmentation must be domain appropriate. Large rotations, severe crops or colour transformations could remove breed-relevant features or create unrealistic examples. Validation and test images must not receive random augmentation because evaluation should be repeatable. The use of modest transformations is consistent with the broader literature on image-data augmentation, which emphasises task validity rather than transformation quantity \citep{shorten2019survey}.

\section{Transfer learning and frozen feature extraction}
Transfer learning reuses knowledge acquired in one task to support another \citep{pan2009survey,yosinski2014transferable}. ImageNet-pretrained CNNs have learned filters from a large and varied image corpus. Although the original output categories differ from restricted-dog classification, low- and mid-level features such as edges, textures, contours and parts remain useful. ``Off-the-shelf'' CNN features have repeatedly provided strong baselines for downstream recognition \citep{sharif2014cnn}.

In frozen feature extraction, all backbone weights are held constant. A new classification head is trained on top. This has several advantages for a small dataset: fewer parameters are updated, training is faster, and the risk of destroying useful general features is reduced. It also supports a fair architecture comparison because each backbone is tested under a common head and training policy. The limitation is that frozen features cannot adapt specifically to the visual boundary between restricted and unrestricted breeds.

\begin{figure}[H]
\centering
\includegraphics[width=0.95\textwidth]{training_pipeline_diagram.pdf}
\caption{Transfer-learning pipeline used conceptually in the project.}
\label{fig:transferpipeline}
\end{figure}

\section{Fine-tuning}
Fine-tuning begins from a trained transfer-learning checkpoint and allows a subset of pretrained layers to update. Upper layers are most closely associated with task-specific combinations, so unfreezing them can adapt the representation while preserving lower-level filters. Fine-tuning is not equivalent to starting again. It is a constrained continuation of training from a strong initial state.

The main risk is catastrophic over-adaptation to the small target dataset. A low learning rate, checkpointing and early stopping are used to reduce this risk. The effectiveness of fine-tuning must be evaluated on held-out data. In the present study, fine-tuning improved precision and F1 while slightly reducing recall. This illustrates why fine-tuning should be treated as an empirical intervention rather than an assumed improvement.

\section{From model score to decision}
The network produces a score, while the evaluation protocol defines a decision. At threshold 0.5, scores greater than or equal to the threshold are classified as restricted. Lowering the threshold generally increases recall but also increases false positives; raising it generally increases precision but can increase false negatives. The selected threshold therefore embodies an operating preference.

Because the dissertation compares frozen and fine-tuned models at a common conventional threshold, the confusion matrices remain directly interpretable. A separate calibration analysis could optimise a threshold for a stated cost function, but it would need to be selected on validation data and evaluated once on untouched test data. Repeatedly selecting a threshold from test performance would leak information and overstate generalisation.


% This chapter can now contain:
% - the ILSVRC development timeline;
% - the revised architecture summary table;
% - the discussion of pre-CNN and CNN milestones.
\chapter{Literature Review and Architecture Comparison}
\label{chap:literature}

\section{ImageNet and the development of deep visual recognition}
Large labelled datasets and accelerated computation changed the dominant approach to image recognition. AlexNet demonstrated that a deep CNN trained with GPUs, ReLU activations, augmentation and dropout could substantially outperform earlier systems on ImageNet \citep{krizhevsky2012imagenet}. Subsequent research increased depth, introduced multi-branch processing, residual connections, separable convolution and automated architecture search. The six backbones evaluated in this dissertation represent several of these major design trajectories.

The historical progression should not be read as a simple ranking in which every newer architecture is superior on every task. Architectural quality interacts with dataset size, input resolution, optimisation, parameter count and transferability. A compact architecture may generalise well and train quickly; a large architecture may offer richer features but impose higher computational cost. Benchmarking under a common pipeline is therefore preferable to selecting a network solely from its original ImageNet result.


\subsection{From hand-crafted features to modern CNN families}
Before the deep-learning breakthrough, large-scale image recognition commonly relied on hand-crafted descriptors and conventional classifiers. AlexNet's 2012 ILSVRC result established the practical advantage of GPU-trained deep CNNs, after which VGG and GoogLeNet/Inception demonstrated alternative routes to greater depth and representational efficiency. ResNet introduced residual connections that made substantially deeper optimisation practical, while later systems increasingly combined multi-scale processing, residual learning, separable convolution and architecture search.

\begin{table}[H]
\centering
\caption{Selected milestones in large-scale visual recognition and their relevance to this study.}
\label{tab:ilsvrctimeline}
\renewcommand{\arraystretch}{1.35}
\begin{tabularx}{0.94\textwidth}{>{\raggedright\arraybackslash}p{0.12\textwidth}>{\raggedright\arraybackslash}p{0.23\textwidth}>{\raggedright\arraybackslash}X}
\toprule
\textbf{Year} & \textbf{Method or architecture} & \textbf{Principal contribution} \\
\midrule
Pre-2012 & Hand-crafted features and shallow classifiers & Feature engineering preceded classification; representations were not learned end to end. \\
2012 & AlexNet & Demonstrated the effectiveness of deep CNNs, GPUs, ReLU and augmentation at ImageNet scale. \\
2014 & VGG / GoogLeNet & Showed the value of uniform small kernels and computationally efficient multi-scale modules. \\
2015 & ResNet & Introduced residual connections for reliable optimisation of very deep networks. \\
2016--2017 & Inception-ResNet / Xception / NASNet & Combined residual and multi-scale ideas, depthwise separable convolution, and architecture search. \\
\bottomrule
\end{tabularx}
\end{table}

The timeline is architectural context rather than a list of the six experimental models or a literal table of annual ILSVRC winners. The project benchmark was selected to compare distinct, established ImageNet-pretrained families available through a common software interface. Consequently, the benchmark contains architectures that were influential without each being the named winner of a particular competition year.

\section{Fine-grained dog-breed recognition}
Dog-breed classification is a well-established fine-grained recognition task. The Stanford Dogs dataset includes substantial intra-class variation: the same breed can appear at different ages, poses, coat lengths, backgrounds and scales. It also exhibits inter-class similarity, as related breeds can share visual characteristics \citep{khosla2011novel}. Fine-grained models must identify discriminative parts while remaining robust to irrelevant context.

The binary reformulation used here differs from conventional 120-class recognition. It reduces output complexity but creates a heterogeneous positive class. The network is not asked to predict a breed name and then map that name to regulation. It learns the binary label directly. This can improve statistical efficiency, but it can also encourage shortcuts. If particular backgrounds, photographic styles or colour distributions correlate with the positive examples, the model may exploit those regularities rather than stable breed morphology. Explanation and error analysis are therefore important complements to aggregate metrics.

\section{VGG16}
VGG16 established the effectiveness of increasing network depth through repeated small $3\times3$ convolutions \citep{simonyan2014very}. The architecture contains thirteen convolutional layers and three fully connected layers in its original ImageNet classifier. Convolutional blocks are separated by max pooling, and channel depth increases as spatial resolution decreases. Its regular structure makes VGG16 easy to explain and inspect.

The main disadvantage is computational inefficiency. The original dense top contributes a very large number of parameters, and even when the top is removed the convolutional backbone remains expensive relative to more modern designs. In transfer learning, VGG features can provide a strong baseline, but the architecture lacks residual connections, branch factorisation and depthwise separability. In this project it achieved the weakest model-zoo result, with 86.48\% accuracy and 58.70\% precision despite recall of 90.53\%. The pattern suggests that it identified many restricted images but generated a comparatively high number of false positives.

\section{ResNet50}
Very deep networks can be difficult to optimise because adding layers may degrade training performance even when overfitting is not the primary issue. ResNet addresses this with identity shortcuts. A residual block learns a residual function $F(x)$ and combines it with the original input:
\begin{equation}
y=F(x)+x.
\end{equation}
The shortcut provides a direct route for information and gradients, making it easier for a block to approximate an identity mapping when additional transformation is unnecessary \citep{he2016deep}.

\begin{figure}[H]
\centering
\includegraphics[width=0.98\textwidth]{residual_block_diagram.pdf}
\caption{Simplified bottleneck residual block. Projection shortcuts may be used when dimensions change.}
\label{fig:residualblock}
\end{figure}

ResNet50 uses bottleneck blocks built from $1\times1$, $3\times3$ and $1\times1$ convolutions. The $1\times1$ layers reduce and then restore channel dimensions, controlling computation. In the model zoo, ResNet50 improved substantially on VGG16 but remained below the Inception-family models, achieving 91.15\% accuracy and an F1 score of 79.55\%.

\section{InceptionV3}
Inception architectures process the same input through several branches and concatenate the resulting feature maps. The branches capture patterns at different effective scales. InceptionV3 improves efficiency through factorisation: a larger convolution can be replaced by sequences of smaller or asymmetric convolutions, such as $1\times7$ followed by $7\times1$ \citep{szegedy2016rethinking}. This reduces computation while preserving a broad receptive field.

The multi-scale design is relevant to dog images because discriminative evidence can be local or global. Ear shape and coat texture are local, whereas body proportion and stance require wider context. InceptionV3 achieved 95.53\% accuracy and 97.37\% recall, indicating that its features transferred effectively to the binary task.

\section{Xception}
Xception interprets an Inception module as an intermediate point between ordinary convolution and depthwise separable convolution \citep{chollet2017xception}. A depthwise separable convolution splits spatial and channel mixing. First, a spatial kernel is applied independently to each channel. A $1\times1$ pointwise convolution then combines information across channels. This factorisation can reduce computation and encourages a distinct representation of spatial and cross-channel relationships.

Xception also uses residual connections. It achieved 94.75\% accuracy and the same 97.37\% recall as InceptionV3, but lower precision. The result demonstrates that high recall alone did not determine selection; InceptionResNetV2 achieved a better overall metric profile.

\section{InceptionResNetV2}
InceptionResNetV2 combines multi-branch Inception processing with residual shortcuts \citep{szegedy2017inceptionresnet}. Its stem is followed by repeated Inception-ResNet-A blocks, a Reduction-A block, repeated Inception-ResNet-B blocks, Reduction-B, and Inception-ResNet-C blocks. Branch outputs are concatenated, projected to the required channel dimension, scaled and added to the shortcut. The original study found that residual connections accelerated the training of Inception networks and could improve performance at similar computational cost.

This architecture offers two plausible advantages for the current task. First, parallel branches represent features at several scales. Second, residual learning supports a very deep network without requiring every block to relearn an entire transformation. Its cost is size and complexity: it has substantially more parameters than NASNetMobile and requires $299\times299$ inputs in the Keras configuration used here.

InceptionResNetV2 produced the strongest frozen result: 97.08\% accuracy, 87.38\% precision, 98.42\% recall, 92.57\% F1, 0.9891 PR-AUC and 0.9101 MCC. The frozen confusion matrix contained 811 true negatives, 27 false positives, 3 false negatives and 187 true positives. This combination of high recall and strong ranking metrics justified its selection for fine-tuning.

\section{NASNetMobile}
NASNet emerged from neural architecture search, in which a controller searches for reusable normal and reduction cells that can be transferred across image scales \citep{zoph2018learning}. NASNetMobile is designed for resource-constrained use. Its relatively low parameter count makes it attractive where deployment size and latency matter.

The architecture achieved 95.14\% accuracy, 80.70\% precision and 96.84\% recall. These results are strong given its compact design, but they did not exceed InceptionResNetV2. For a mobile deployment study, computational efficiency might receive greater weight. In this dissertation, predictive performance and controlled comparison were prioritised over on-device inference.

\section{Architectural summary}
Table~\ref{tab:architectureoverview} summarises the design ideas. ``Depth'' is not perfectly standardised across publications and software libraries: paper names such as ResNet50 refer to conventional weighted-layer counts, while Keras documentation can report implementation-level layer depth. The table therefore emphasises structural families and parameter scale rather than treating depth as a directly comparable physical measurement.

\begin{table}[H]
\centering
\caption{Architectural overview of the six pretrained backbones. Parameter counts are approximate Keras application totals with the original classification top.}
\label{tab:architectureoverview}
\footnotesize
\renewcommand{\arraystretch}{1.45}
\setlength{\tabcolsep}{3.5pt}
\begin{tabularx}{0.98\textwidth}{>{\raggedright\arraybackslash}p{2.45cm} >{\centering\arraybackslash}p{1.35cm} >{\centering\arraybackslash}p{1.45cm} >{\raggedright\arraybackslash}X >{\raggedright\arraybackslash}X}
\toprule
Architecture & Input & Parameters & Core design & Main implication \\
\midrule
VGG16 & $224^2$ & 138.4 M & Repeated $3\times3$ convolution and max-pooling blocks & Simple and interpretable, but parameter-heavy \\
ResNet50 & $224^2$ & 25.6 M & Bottleneck residual blocks with identity shortcuts & Enables deeper optimisation and feature reuse \\
InceptionV3 & $299^2$ & 23.9 M & Parallel multi-scale branches and factorised convolution & Efficient representation of features at several scales \\
Xception & $299^2$ & 22.9 M & Depthwise separable convolution with residual flow & Separates spatial filtering from channel mixing \\
InceptionResNetV2 & $299^2$ & 55.9 M & Inception branches combined with residual shortcuts & High-capacity multi-scale representation \\
NASNetMobile & $224^2$ & 5.3 M & Searched normal and reduction cells & Compact architecture suited to limited-resource inference \\
\bottomrule
\end{tabularx}
\end{table}

\section{Transfer learning literature}
Transfer learning is motivated by the observation that learned representations vary in their specificity. Lower layers tend to be more general, while later layers become increasingly specialised to the source task \citep{yosinski2014transferable}. Reusing a pretrained backbone can therefore reduce data requirements and training time. The success of frozen CNN features as general-purpose descriptors was demonstrated even before transfer learning became routine in high-level libraries \citep{sharif2014cnn}.

Two strategies are common. Feature extraction freezes the backbone and learns a new head. Fine-tuning subsequently updates part of the backbone. The second strategy can improve domain alignment, but it also increases variance and overfitting risk. A rigorous study must therefore compare held-out results rather than assuming that more trainable parameters produce a better model.

\section{Evaluation under imbalance}
Accuracy is the proportion of correct predictions, but it can obscure minority-class behaviour. Precision measures the proportion of positive predictions that are correct, while recall measures the proportion of actual positives recovered. F1 is their harmonic mean. ROC-AUC evaluates ranking across false-positive rates, whereas PR-AUC focuses directly on positive predictive performance and is often more informative when the positive class is less prevalent \citep{davis2006relationship}. Matthews correlation coefficient summarises all four confusion-matrix cells and remains informative under imbalance \citep{matthews1975comparison}.

No single metric defines the optimal model independently of use. A screening system may prioritise recall; a system in which false accusations are particularly harmful may prioritise precision. The present selection considered recall, PR-AUC, F1 and MCC together. This avoids hiding an unacceptable error type behind a high aggregate score.

\section{Interpretability and Grad-CAM}
Grad-CAM computes gradients of a target output with respect to feature maps in a convolutional layer, averages those gradients to obtain channel weights, and combines the weighted maps into a coarse localisation heatmap \citep{selvaraju2017gradcam}. It can indicate which spatial regions are associated with a decision without retraining the model.

The method has important limitations. A visually plausible heatmap does not prove that the model used the same reasoning as a human. Saliency methods can fail sensitivity tests, and visually attractive explanations can remain similar after model or label randomisation \citep{adebayo2018sanity}. Grad-CAM is spatially coarse because it is based on late convolutional maps. It is best used as a diagnostic aid rather than a legal explanation.

In high-stakes settings, post-hoc explanation may be insufficient if the underlying model is opaque and contestability is weak \citep{rudin2019stop}. The present project uses Grad-CAM to interrogate a research artefact, not to legitimise autonomous deployment.

\section{Fairness, documentation and high-stakes use}
Bias in image classification can arise from sampling, labels, camera conditions, background correlations and uneven representation. Fairness is not limited to demographic attributes; here it also concerns unequal model reliability across breeds, body types, age groups and photographic conditions. Dataset documentation and model cards provide structured ways to state intended use, composition, limitations and evaluation conditions \citep{gebru2021datasheets,mitchell2019modelcards}.

A legally consequential classifier would require more than accuracy. It would require validated definitions, representative field data, calibration, uncertainty estimates, human review, appeal mechanisms, audit trails and monitoring after deployment. The literature therefore supports a cautious distinction between technical feasibility and institutional legitimacy.

\section{Architecture depth, parameters and computational cost}
Architecture comparison is often reduced to parameter count, but parameters are only one source of cost. Convolutional computation also depends on input resolution, feature-map dimensions, channel counts and kernel operations. InceptionResNetV2 has more parameters than the other modern backbones in this study and processes $299\times299$ inputs. It therefore imposes a greater memory and inference burden than NASNetMobile. The model was selected because predictive performance was the primary experimental objective, not because it was the most efficient option.

Parameter count also has different meanings depending on whether the original ImageNet top is included. The Keras application totals commonly quoted for VGG16 include its large dense classifier. In this project, that top was removed and replaced by global average pooling and a binary head. Consequently, the number of trainable parameters in the frozen stage was much smaller than the architecture's published total. The pretrained backbone parameters remained present in memory but did not receive gradient updates.

Depth is similarly ambiguous. ResNet50 conventionally names a 50-layer weighted architecture. Inception networks contain branches and auxiliary operations that software libraries can count in several ways. A table that compares ``layers'' without defining the counting rule can imply false precision. For this reason, the report uses conventional names, broad Keras depth information and block-level descriptions rather than asserting that one implementation-level layer count is directly comparable with another.

\section{Why EfficientNet and vision transformers were not included}
EfficientNet and vision transformers are important modern alternatives, but the final model zoo was intentionally bounded. EfficientNet scales depth, width and resolution through a compound rule, while vision transformers divide an image into patches and learn attention-based relationships. Including them would broaden the historical comparison, but it would also increase experimental multiplicity and require additional decisions about augmentation, regularisation and pretrained checkpoints.

The six selected CNNs already span major architectural ideas and answer the principal question within the available compute. Excluding newer families is therefore a scope decision, not a claim that they are inferior. Future work could compare EfficientNetV2, ConvNeXt and transformer-based backbones under the same split. Such a comparison should predeclare the model set and training budget to avoid selecting architectures only after favourable results are observed.

\section{Dog-breed identification and human reliability}
A key background issue is that visual breed labels are not perfectly objective, especially for mixed dogs. Studies of shelter identification have reported disagreement between visual assessments and DNA-derived breed assignments. This matters because supervised learning inherits the certainty and uncertainty of its labels. Even a model that reproduces dataset labels with high accuracy may not reproduce genetic ancestry or legal interpretation in ambiguous cases.

Purebred benchmark datasets reduce this ambiguity by presenting named categories, which helps controlled experimentation. The same feature is also a limitation: the benchmark may make the problem substantially easier than field identification. The gap between benchmark labels and operational definitions is therefore part of the validity argument, not a reason to discard the experiment. The model can be evaluated accurately for the task it was trained to perform while remaining unsuitable for a broader task.

\section{Literature-review synthesis}
The literature supports four design choices in this study. First, CNNs remain strong visual feature extractors because local connectivity, hierarchical representation and weight sharing match image structure. Second, transfer learning is appropriate where target data are limited. Third, modern architectures embody different efficiency and optimisation strategies, making controlled comparison preferable to arbitrary selection. Fourth, high-stakes use requires explanation, documentation and governance beyond conventional test accuracy.

The literature also identifies unresolved tensions. Greater model capacity can improve benchmark performance while reducing transparency and increasing compute. Post-hoc heatmaps can aid inspection while creating an illusion of explanation. A binary legal label can simplify modelling while obscuring heterogeneous breed-level behaviour. These tensions frame the interpretation of the experimental results.


% This chapter can now contain:
% - the three-stage experimental pipeline;
% - the dataset summary table;
% - the notebook and reproducibility workflow.
\chapter{Methodology}
\label{chap:methodology}

\section{Research design}
The project used a staged experimental methodology. The independent experimental factor was the pretrained CNN architecture and, later, training stage. The dependent outcomes were test-set classification metrics and confusion-matrix counts. Dataset composition, split, class order, evaluation code and general head design were held constant as far as architecture-specific input requirements allowed.

The workflow comprised four notebooks: dataset preparation; frozen model-zoo benchmarking; fine-tuning of the selected model; and Grad-CAM explanation. Each notebook saved outputs to persistent storage. This separation reduced the risk that an interrupted Colab runtime would erase results and created a traceable sequence of artefacts.

\section{Experiment at a glance}
The dissertation used a controlled two-stage transfer-learning experiment. First, six ImageNet-pretrained CNN backbones were frozen and evaluated under the same dataset, preprocessing, classification head, training controls and held-out test set. Second, the strongest frozen architecture was selected using metric-aware criteria and its final 20 layers were unfrozen for limited fine-tuning. This design separates architecture comparison from subsequent task-specific adaptation.

\begin{table}[H]
\centering
\caption{Summary of the experimental pipeline.}
\label{tab:experimentglance}
\renewcommand{\arraystretch}{1.4}
\begin{tabularx}{0.96\textwidth}{p{0.25\textwidth}X}
\toprule
\textbf{Stage} & \textbf{Implementation} \\
\midrule
Data construction & Stanford Dogs images; six represented restricted breeds and 24 unrestricted breeds selected reproducibly. \\
Split design & Fixed training, validation and held-out test partitions reused across all models. \\
Frozen benchmark & VGG16, ResNet50, InceptionV3, Xception, InceptionResNetV2 and NASNetMobile trained as frozen feature extractors. \\
Model selection & Multiple imbalance-aware metrics and confusion counts considered rather than accuracy alone. \\
Fine-tuning & Final 20 layers of the selected InceptionResNetV2 backbone unfrozen at a low learning rate. \\
Explanation & Fixed true-positive, true-negative, false-positive and false-negative cases analysed using Grad-CAM. \\
Reproducibility & Notebooks, saved figures, tables, predictions, checkpoints and LaTeX source retained in the project archive. \\
\bottomrule
\end{tabularx}
\end{table}

\begin{figure}[H]
\centering
\includegraphics[width=0.96\textwidth]{training_pipeline_diagram.pdf}
\caption{End-to-end experimental pipeline from reproducible dataset construction to frozen benchmarking, selected-model fine-tuning, evaluation and explanation.}
\label{fig:trainingpipeline}
\end{figure}

\section{Dataset source and class construction}
Stanford Dogs contains 120 breed directories and more than twenty thousand images \citep{khosla2011novel}. The preparation notebook verified the extracted directory structure before constructing the binary dataset. Six breeds represented both in the dataset and the restricted list were assigned to the positive class: Bull Mastiff, Doberman Pinscher, German Shepherd, Rhodesian Ridgeback, Rottweiler and Staffordshire Bull Terrier.

Twenty-four unrestricted breeds were selected using a fixed random seed. Fixing this selection is important. If unrestricted breeds were re-sampled for every run, model comparisons could reflect a change in negative-class difficulty rather than architectural performance. The prepared split contained 3,023 training images, 1,002 validation images and 1,028 test images. Table~\ref{tab:datasetcounts} shows the class distribution.

\begin{table}[H]
\centering
\caption{Prepared dataset counts.}
\label{tab:datasetcounts}
\renewcommand{\arraystretch}{1.6}
\begin{tabular}{lrrr}
\toprule
Split & Unrestricted & Restricted & Total \\
\midrule
Training & 2,462 & 561 & 3,023 \\
Validation & 816 & 186 & 1,002 \\
Test & 838 & 190 & 1,028 \\
\bottomrule
\end{tabular}
\end{table}

The positive prevalence in the test set was approximately 18.48\%. This imbalance is moderate but sufficient to make accuracy alone misleading. A classifier predicting every image as unrestricted would achieve about 81.52\% accuracy while having zero restricted-class recall.

\section{Split integrity and leakage control}
The data were divided into training, validation and test sets using a 60/20/20 strategy at the image level. The training set was used for parameter updates, the validation set for checkpoint and stopping decisions, and the test set for final comparison. The test set was not intended as an iterative tuning resource.

A limitation is that the source dataset can contain visually related or near-duplicate web images. The project did not implement perceptual-hash duplicate detection across splits. Consequently, the reported performance should be interpreted as performance on the prepared Stanford-derived split, not as a guarantee of independence at the level of original photographic events. Future work should audit exact and near duplicates before splitting.

\section{Preprocessing and augmentation}
Each architecture was paired with its expected input resolution and preprocessing function. InceptionV3, Xception and InceptionResNetV2 used $299\times299$ inputs, while VGG16, ResNet50 and NASNetMobile used $224\times224$ inputs in their standard Keras configurations. Images were decoded as RGB and batched through the TensorFlow data pipeline.

The training graph included light horizontal flipping, rotation and zoom. Validation and test data were preprocessed deterministically without random augmentation. This separation ensures that repeated evaluation of the same checkpoint produces the same predictions. Augmentation was deliberately modest because strong geometric transformation could obscure breed morphology.

\section{Model-zoo configuration}
Each model used ImageNet-pretrained weights with the original classification top removed. The backbone was frozen. Its output was passed through global average pooling, dropout and a one-unit sigmoid classifier. This standardised head allowed the comparison to focus on the transferred representation rather than a separately engineered classifier for each model.

The architectures were VGG16, ResNet50, InceptionV3, Xception, InceptionResNetV2 and NASNetMobile. These models span conventional deep convolution, residual learning, multi-branch factorisation, depthwise separable convolution, hybrid Inception-residual design and neural architecture search. The benchmark therefore tests qualitatively different feature extractors rather than six minor variants of one family.

Class weighting was used to reduce majority-class dominance during optimisation. The weights alter the loss contribution of examples, not the reported test counts. They do not fabricate balance in the test set; they encourage the optimiser to treat errors on the less frequent class as more consequential during training.

\section{Training controls}
The head was optimised with binary cross-entropy and Adam. Checkpointing preserved the model with the best monitored validation result, and early stopping limited continued training after improvement ceased. Training histories and model files were saved under architecture-specific names.

A fair benchmark does not require every architecture to experience identical internal computation, because their input sizes and normalisation differ. It does require the same dataset split, target labels, evaluation procedure and broad training policy. The model-zoo notebook encoded these controls in one loop to reduce manual inconsistency.

\section{Metric definitions}
Let $TP$, $TN$, $FP$ and $FN$ denote true positives, true negatives, false positives and false negatives. The primary metrics are
\begin{align}
\mathrm{Accuracy}&=\frac{TP+TN}{TP+TN+FP+FN},\\
\mathrm{Precision}&=\frac{TP}{TP+FP},\\
\mathrm{Recall}&=\frac{TP}{TP+FN},\\
F_1&=2\frac{\mathrm{Precision}\times\mathrm{Recall}}{\mathrm{Precision}+\mathrm{Recall}}.
\end{align}

Specificity was also available from $TN/(TN+FP)$. ROC-AUC evaluates the ranking of positive and negative scores across thresholds. PR-AUC summarises the precision--recall curve and directly reflects performance on the positive class. MCC is
\begin{equation}
\mathrm{MCC}=\frac{TP\times TN-FP\times FN}{\sqrt{(TP+FP)(TP+FN)(TN+FP)(TN+FN)}}.
\end{equation}
MCC approaches one for perfect classification, zero for no better than chance association, and minus one for complete disagreement.

\section{Model selection}
The strongest frozen model was selected through metric-aware ranking rather than accuracy alone. High restricted-class recall was important, but precision, F1, PR-AUC and MCC were also considered. InceptionResNetV2 ranked strongest overall and was therefore advanced to fine-tuning.

This selection occurred before examination of the fine-tuned test result. The process avoids the less defensible practice of fine-tuning every model and then reporting only the most favourable outcome without a clear selection rule. It also limits computational cost.

\section{Fine-tuning strategy}
The best frozen InceptionResNetV2 checkpoint was loaded. The final twenty layers of the pretrained base were made trainable, while the remainder stayed frozen. The model was recompiled with a learning rate of $1\times10^{-5}$. Recompilation is essential after changing trainable status because the optimiser must construct its parameter update set from the new configuration.

Fine-tuning was controlled by validation monitoring and checkpointing. The final reported checkpoint was evaluated on the same held-out test set as the frozen model. This permits a direct comparison but also means the two evaluations are statistically paired: they classify the same images. The report focuses on observed differences rather than claiming formal significance from one test split.

\section{Explanation sample}
The explanation notebook generated predictions for the complete test set and assigned each case to TP, TN, FP or FN. Three cases from each category were selected according to a fixed procedure, creating twelve examples. Grad-CAM heatmaps were computed from the final convolutional feature layer of the fine-tuned InceptionResNetV2 model and overlaid on the original image.

The explanation sample is qualitative. It was designed to expose varied outcomes, not to estimate the prevalence of particular attention patterns. Selecting equal numbers from each confusion category intentionally over-represents errors relative to their frequency in the full test set. This is appropriate for diagnostic analysis but must not be mistaken for a random sample of model behaviour.

\section{Reproducibility and project management}
Random seeds were set where supported. Dataset folders, class order and split counts were verified before training. Results were exported as CSV files, histories were saved, and confusion matrices and learning curves were written as standalone images. The report archive contains the four principal notebooks and the figures used in the dissertation.

The staged workflow also reflects iterative project management. The original proposal considered a broader architecture timeline and hyperparameter search. Resource constraints and the need for a clear, reproducible final experiment led to a focused six-model benchmark followed by one controlled fine-tuning stage. This was a methodological narrowing, not an unreported failure: the completed design addresses the research questions with less uncontrolled multiplicity.

\section{Ethical methodology}
The source data contain dog photographs and breed labels rather than human personal data, but ethical risk remains. The central risk is downstream misuse: a model score could be treated as objective proof in a context where visual breed identification is contested. The methodology therefore treats legal status as a dataset label and avoids claims about aggression or individual dangerousness.

No attempt was made to infer temperament. The model cannot assess behaviour, training, ownership conditions or context. The report explicitly separates classification performance from policy validity. This distinction is necessary because a technically successful classifier could still support an unjust or scientifically disputed decision framework.

\section{Software environment}
The experiments were implemented in Python using TensorFlow and Keras within Google Colab. Google Drive provided persistent storage for data splits, checkpoints, prediction files and plots. GPU acceleration was used for model-zoo training and fine-tuning. The notebooks verify whether the expected directories and hardware resources are available before beginning expensive stages.

The use of a managed notebook environment improves accessibility but introduces reproducibility risks. Library versions can change, runtimes are ephemeral and mounted-drive paths can become inconsistent across revisions. The final notebooks therefore record TensorFlow versions in their outputs and use explicit project paths. A stronger archival package would additionally include a locked requirements file or container image. The report treats the notebooks and saved artefacts as the reproducibility record for the completed run rather than promising bit-for-bit equivalence on future software versions.

\section{Data pipeline performance considerations}
Image training can become input-bound if data are repeatedly decoded from slow remote storage. The preparation notebook extracted and organised files with attention to local versus mounted-drive performance. The TensorFlow pipeline used batching and prefetching so that CPU-side image preparation could overlap with GPU computation where possible. These engineering choices do not change the statistical model, but they affect whether the experiment is feasible within a hosted runtime.

The persistent split directories also reduce accidental changes between notebooks. Model benchmarking, fine-tuning and explanation all point to the same train, validation and test folders. Rebuilding the split is unnecessary unless the data-preparation stage intentionally changes. The printed verification counts provide a simple integrity check before training.

\section{Decision not to perform exhaustive hyperparameter search}
The proposal considered systematic search over dropout, learning rate and other settings. The final study used a restrained configuration instead. Exhaustive search would consume substantial GPU time and, without a nested validation design, could overfit the validation set. It would also make the architecture comparison less interpretable if each network received a different and unequally extensive search.

A common pipeline provides a cleaner answer to the comparative question. It may not reveal the maximum possible performance of every model, but it estimates how their pretrained features behave under a consistent practical setup. The selected model then receives one clearly described fine-tuning intervention. This design prioritises methodological clarity over leaderboard optimisation.

\section{Threats to methodological validity}
Internal validity could be weakened by hidden duplicates, inconsistent preprocessing or accidental test use. The notebook structure and common evaluation functions reduce but do not eliminate these risks. Construct validity is limited because the binary label operationalises legal membership using dataset categories rather than direct statutory adjudication. External validity is limited by the curated source images. Conclusion validity is limited by one split and one main run per architecture.

Stating these threats is not an admission that the experiment is invalid. It defines the level of inference supported. The strongest claim is comparative: under this prepared split and pipeline, InceptionResNetV2 produced the strongest frozen metric profile and fine-tuning altered the error balance in the observed way.

\section{Research integrity and reporting choices}
The report preserves unfavourable aspects of the results. It reports the reduction in recall after fine-tuning, discusses background attention, and avoids presenting Grad-CAM only for successful cases. It also distinguishes the established 0.5-threshold confusion matrix from exploratory calibration output. These choices reduce selective reporting.

The experimental figures are used directly from saved outputs rather than redrawn with altered values. Numerical tables are rounded for readability, while precise values remain in the text and source CSV files. Rounding to two decimal places in summary tables avoids an impression of precision beyond the finite test set.


% This chapter can now contain:
% - the revised frozen benchmark table;
% - the regenerated four-metric bar chart;
% - the frozen/fine-tuned summary table;
% - confusion matrices and learning curves.
\chapter{Results}
\label{chap:results}

\section{Frozen model-zoo comparison}
Table~\ref{tab:modelzoo} presents the principal frozen-model metrics. Values are shown to two decimal places as percentages in the narrative and to two decimal places on the zero-to-one scale in the table. InceptionResNetV2 achieved the strongest overall result. VGG16 was clearly weakest, while the remaining four modern architectures formed a stronger group.

\begin{table}[H]
\centering
\caption{Frozen model-zoo test performance.}
\label{tab:modelzoo}
\small
\renewcommand{\arraystretch}{1.85}
\setlength{\tabcolsep}{4.5pt}
\begin{tabularx}{\textwidth}{>{\raggedright\arraybackslash}X c c c c c}
\toprule
Model & Accuracy & Precision & Recall & F1 & PR-AUC \\
\midrule
VGG16 & 0.86 & 0.59 & 0.91 & 0.71 & 0.81 \\
ResNet50 & 0.91 & 0.69 & 0.93 & 0.80 & 0.89 \\
InceptionV3 & 0.96 & 0.82 & 0.97 & 0.89 & 0.96 \\
Xception & 0.95 & 0.79 & 0.97 & 0.87 & 0.97 \\
InceptionResNetV2 & \textbf{0.97} & \textbf{0.87} & \textbf{0.98} & \textbf{0.93} & \textbf{0.99} \\
NASNetMobile & 0.95 & 0.81 & 0.97 & 0.88 & 0.96 \\
\bottomrule
\end{tabularx}
\end{table}

\begin{figure}[H]
\centering
\includegraphics[width=0.96\textwidth]{model_zoo_four_metric_comparison_presentation.pdf}
\caption{Comparison of accuracy, precision, recall and F1 across the frozen model zoo.}
\label{fig:modelzoocomparison}
\end{figure}

The contrast between VGG16 and the Inception-family models is especially visible in precision. VGG16 recovered 90.53\% of restricted images, but only 58.70\% of its restricted predictions were correct. This suggests a broad positive decision region that generated many false alarms. ResNet50 improved precision to 69.41\%, while InceptionV3, Xception, NASNetMobile and InceptionResNetV2 exceeded 79\%.

InceptionV3 and Xception both achieved recall of 97.37\%. Their difference lay primarily in precision and ranking metrics. InceptionV3 reached 81.86\% precision and 88.94\% F1; Xception reached 79.06\% precision and 87.26\% F1. NASNetMobile reached 80.70\% precision, 96.84\% recall and 88.04\% F1. The compact NASNet model was therefore competitive, although not the strongest.


\subsection{Benchmark decision}
The benchmark identified InceptionResNetV2 as the strongest frozen model overall. Its principal advantage was not accuracy alone: it combined the highest recall and F1 with strong precision, PR-AUC and MCC. The model was therefore advanced to fine-tuning under the selection rule defined in Chapter~\ref{chap:methodology}.

\section{Selected frozen model}
InceptionResNetV2 achieved 97.08\% accuracy, 87.38\% precision, 98.42\% recall, 92.57\% F1, 0.9974 ROC-AUC, 0.9891 PR-AUC and 0.9101 MCC. Its confusion matrix contained 811 true negatives, 27 false positives, 3 false negatives and 187 true positives.

\begin{figure}[H]
\centering
\includegraphics[width=0.62\textwidth]{InceptionResNetV2_confusion_matrix.png}
\caption{Confusion matrix for frozen InceptionResNetV2.}
\label{fig:frozenconfusion}
\end{figure}

The model missed only three of 190 restricted images. This is a recall of $187/190=98.42\%$. However, 27 unrestricted images were labelled restricted. Specificity was $811/838=96.78\%$. The result is excellent as a controlled benchmark, but the absolute counts matter: a high percentage does not erase the practical meaning of thirty incorrect decisions.

\section{Fine-tuned performance}
Fine-tuning improved accuracy to 98.25\%, precision to 93.88\%, F1 to 95.34\% and MCC to 0.9428. Recall declined to 96.84\%. The fine-tuned confusion matrix contained 826 true negatives, 12 false positives, 6 false negatives and 184 true positives.

\begin{table}[H]
\centering
\caption{Frozen and fine-tuned InceptionResNetV2 comparison.}
\label{tab:frozenfine}
\renewcommand{\arraystretch}{1.85}
\setlength{\tabcolsep}{7pt}
\begin{tabular}{lrrrr}
\toprule
Stage & Accuracy & Precision & Recall & F1 \\
\midrule
Frozen & 0.97 & 0.87 & \textbf{0.98} & 0.93 \\
Fine-tuned & \textbf{0.98} & \textbf{0.94} & 0.97 & \textbf{0.95} \\
\bottomrule
\end{tabular}
\end{table}

\begin{quote}
\textbf{Central result.} Fine-tuning improved overall balance and positive predictive reliability, but slightly reduced sensitivity to restricted examples. False positives fell by 15, from 27 to 12, while false negatives increased by 3, from 3 to 6.
\end{quote}

\begin{figure}[H]
\centering
\includegraphics[width=0.62\textwidth]{FineTuned_InceptionResNetV2_confusion_matrix.png}
\caption{Confusion matrix for fine-tuned InceptionResNetV2.}
\label{fig:fineconfusion}
\end{figure}

The most important change is not the 1.17 percentage-point increase in accuracy. False positives fell from 27 to 12, a reduction of 15 cases or 55.56\%. False negatives increased from 3 to 6. Fine-tuning therefore made the classifier more conservative about assigning the restricted label. This improved precision by 6.50 percentage points but reduced recall by 1.58 points.

Specificity increased from 96.78\% to 98.57\%. The positive predictive value increased because fewer unrestricted images were included among positive predictions. Whether this is preferable depends on the cost model. If missing any restricted case is considered paramount, the frozen model's three false negatives may be preferred. If false designation of unrestricted dogs is considered especially harmful, the fine-tuned model provides a better balance.

\section{Learning behaviour}
\begin{figure}[H]
\centering
\includegraphics[width=0.88\textwidth]{FineTuned_InceptionResNetV2_learning_curves.png}
\caption{Training and validation curves for fine-tuned InceptionResNetV2.}
\label{fig:finecurves}
\end{figure}

The curves show the optimisation trajectory used for checkpoint selection. Training performance alone is not the selection criterion because a high-capacity model can continue to improve on training examples while validation loss worsens. The saved checkpoint and early-stopping controls protect against simply reporting the final epoch.

The curves should also be interpreted cautiously. Validation performance is a noisy estimate based on 1,002 images, and small fluctuations can alter the nominal best epoch. A repeated-seed study or cross-validation would provide stronger estimates of variance. The available results demonstrate one reproducible run rather than a confidence interval over possible training histories.

\section{Metric consistency}
The fine-tuned model's higher accuracy, precision, F1 and MCC are mutually consistent with the confusion matrix. Its ROC-AUC and PR-AUC were marginally lower than the frozen model in the saved comparison, despite better thresholded performance. This is possible because AUC measures ranking across thresholds, whereas confusion-matrix metrics describe one threshold. Fine-tuning improved the chosen operating point without necessarily improving score ordering everywhere.

This distinction is central to responsible reporting. Saying that the fine-tuned model is ``better'' without qualification would hide the loss in recall and the small AUC changes. The defensible conclusion is that fine-tuning produced a more balanced thresholded classifier with substantially fewer false positives, while the frozen model retained the highest sensitivity to restricted examples.

\section{Error concentration}
All three selected false-negative explanation cases were Staffordshire Bull Terriers. This does not prove that every false negative in the complete model is from that breed without a full per-breed table, but it identifies an important diagnostic pattern in the fixed sample. Staffordshire Bull Terriers display substantial variation in pose, expression, body condition and photographic framing, and they can share features with both restricted and unrestricted bull-type breeds.

The false-positive examples included a Black-and-Tan Coonhound, a Kelpie and a Weimaraner. These cases demonstrate different routes to error: colour and head pattern, shepherd-like silhouette and contextual or compositional cues. The next chapter examines these examples through Grad-CAM.

\section{Relative improvement and practical interpretation}
The change from 97.08\% to 98.25\% accuracy is a relative reduction in total error from 30 incorrect cases to 18. This is a 40\% reduction in the number of test errors. The distribution of that reduction matters: fifteen false positives were removed while three additional false negatives were introduced. The net change of twelve fewer errors therefore conceals movement in opposite directions.

Precision rose because the denominator $TP+FP$ became substantially cleaner. The frozen model made 214 positive predictions, of which 187 were correct. The fine-tuned model made 196 positive predictions, of which 184 were correct. It generated fewer positive alerts and a larger fraction of those alerts were correct. Recall declined because three of the previously recovered positives moved below the threshold.

The F1 increase from 92.57\% to 95.34\% reflects the stronger harmonic balance. MCC also increased, indicating improvement across all four confusion cells rather than only the majority class. These measures support selection of the fine-tuned model as the more balanced final classifier, provided the stated preference does not demand maximum possible recall.

\section{Why AUC and threshold metrics can disagree}
AUC statistics consider all possible thresholds and are based on score ordering. A confusion matrix considers one threshold. Fine-tuning can move scores so that the 0.5 boundary becomes better aligned with the classes while slightly worsening the global ordering of some pairs. This explains why thresholded precision and F1 can improve even when ROC-AUC or PR-AUC changes marginally in the other direction.

For deployment, both views matter. AUC indicates whether a useful operating point may exist; threshold metrics describe the selected operating point. A model with excellent AUC still requires a justified threshold. Conversely, a strong confusion matrix at one threshold can mask poor calibration or instability when prevalence changes.

\section{Class prevalence and predictive value}
Precision depends on prevalence. The test prevalence of restricted images is about 18.48\%. In a population with a much lower prevalence, the same sensitivity and specificity could produce a lower positive predictive value because false positives would arise from a larger negative pool. This is the base-rate effect.

Therefore, the reported 93.88\% precision cannot be transferred directly to an operational population without knowing its case selection process and prevalence. If the model were applied only after expert triage, prevalence could be higher; if applied indiscriminately to general dog photographs, prevalence could be much lower. This is another reason to avoid interpreting benchmark precision as a universal property.

\section{Results in relation to the hypotheses}
The results support the hypothesis that pretrained deep architectures can distinguish the constructed classes effectively. They also support the expectation that architectures differ materially under common conditions. The result for InceptionResNetV2 is consistent with the hypothesis that modern multi-scale residual features would transfer strongly.

The fine-tuning hypothesis is only partly supported. Fine-tuning improved several headline metrics, but it did not improve recall or every ranking measure. The correct conclusion is conditional: fine-tuning improved the balance of the chosen threshold, not every dimension of performance. This partial confirmation is more informative than forcing the outcome into a simple success/failure statement.


\chapter{Model Explanation with Grad-CAM}
\label{chap:explanation}

\section{Purpose and scope}

The quantitative evaluation established that the selected InceptionResNetV2 classifier performed strongly on the controlled test set. However, metrics such as accuracy, precision, recall, F1 score, PR-AUC, MCC, and the confusion matrix describe the correctness of the predictions rather than the visual evidence used to produce them. Gradient-weighted Class Activation Mapping (Grad-CAM) was therefore applied as a post-hoc explanation technique \citep{selvaraju2017gradcam}.

Grad-CAM uses the gradients of a target class flowing into a late convolutional representation to produce a coarse class-discriminative localisation map. The resulting heatmap is resized and overlaid on the original image. Warmer colours represent regions with comparatively stronger influence on the selected class score. The method is useful for inspecting whether a model attends to the dog, to plausible breed-related morphology, or to incidental contextual features. It is not a pixel-level segmentation method, and a visually plausible heatmap does not prove that the classifier has learned a robust or causal concept \citep{adebayo2018sanity,rudin2019stop}.

\section{Selection and interpretation of cases}

A reproducible sample of twelve test images was selected: three true positives, three true negatives, three false positives, and three false negatives. Each saved figure contains the original image, the Grad-CAM heatmap, and the merged overlay. The original image was considered alongside the heatmap because activation colours are meaningful only in relation to the visual content of the scene.

The examples were examined using human visual judgement. Particular attention was paid to the apparent size of the dog in the frame, pose, occlusion, facial expression, contrast, background texture, and whether the strongest activation overlapped the dog or an unrelated object. This is a qualitative analysis intended to complement, rather than replace, the principal metric-based evaluation.

\section{False negatives: restricted dogs predicted as unrestricted}

\begin{figure}[p]
    \centering
    \begin{subfigure}[t]{0.84\textwidth}
        \centering
        \includegraphics[width=\linewidth]{FN_staffordshire_bullterrier_n02093256_2987_gradcam.png}
        \caption{Close portrait with an unusual facial expression and mouth position. Activation is broad rather than tightly focused on the muzzle, skull, or body proportions.}
        \label{fig:fn-portrait}
    \end{subfigure}

    \vspace{0.8em}
    \begin{subfigure}[t]{0.84\textwidth}
        \centering
        \includegraphics[width=\linewidth]{FN_staffordshire_bullterrier_n02093256_5325_gradcam.png}
        \caption{A strong high-contrast structure competes with the dog for attention, producing a response distributed across both subject and background.}
        \label{fig:fn-pattern}
    \end{subfigure}

    \vspace{0.8em}
    \begin{subfigure}[t]{0.84\textwidth}
        \centering
        \includegraphics[width=\linewidth]{FN_staffordshire_bullterrier_n02093256_5988_gradcam.png}
        \caption{The standing dog occupies a relatively small part of a visually active frame, and the activation remains diffuse.}
        \label{fig:fn-standing}
    \end{subfigure}
    \caption{Representative false-negative Grad-CAM explanations. All three restricted Staffordshire Bull Terriers were classified as unrestricted.}
    \label{fig:gradcam-fn}
\end{figure}

The false negatives are visually understandable. They contain difficult composition, competing contrast, unusual expression, or a dog that is not presented as a clean reference image. In each case, the heatmap is less convincingly concentrated on the most recognisable parts of the animal. The examples suggest that restricted-class sensitivity can fall when the animal is relatively small, visually unusual, or embedded in a strong background.

From a human perspective, this is consistent with the nature of fine-grained classification. A person also depends on a sufficiently clear view of the head, muzzle, body proportions, coat, and stance. If those features are distorted, partly hidden, or visually overwhelmed by the scene, the distinction becomes more difficult.

\section{False positives: unrestricted dogs predicted as restricted}

\begin{figure}[p]
    \centering
    \begin{subfigure}[t]{0.84\textwidth}
        \centering
        \includegraphics[width=\linewidth]{FP_black-and-tan_coonhound_n02089078_1454_gradcam.png}
        \caption{The Black and Tan Coonhound was predicted as restricted. Strong activation around the head and upper body suggests genuine morphological similarity rather than background alone.}
        \label{fig:fp-coonhound}
    \end{subfigure}

    \vspace{0.8em}
    \begin{subfigure}[t]{0.84\textwidth}
        \centering
        \includegraphics[width=\linewidth]{FP_kelpie_n02105412_7370_gradcam.png}
        \caption{The Kelpie was predicted as restricted with high confidence. The heatmap overlaps the dog but also extends into the street environment.}
        \label{fig:fp-kelpie}
    \end{subfigure}

    \vspace{0.8em}
    \begin{subfigure}[t]{0.84\textwidth}
        \centering
        \includegraphics[width=\linewidth]{FP_weimaraner_n02092339_3028_gradcam.png}
        \caption{The Weimaraner was classified as restricted. Activation covers both the dog and the highly textured flower background.}
        \label{fig:fp-weimaraner}
    \end{subfigure}
    \caption{Representative false-positive Grad-CAM explanations. Unrestricted dogs were incorrectly classified as restricted.}
    \label{fig:gradcam-fp}
\end{figure}

The false positives show why the classifier should not be interpreted as a legal breed-identification system. Some unrestricted breeds share visible characteristics with breeds in the restricted class, including head shape, body build, coat colour, stance, or silhouette. The heatmaps frequently overlap the dog, so the errors cannot be dismissed simply as background failure. Instead, the learned decision boundary appears to treat some legitimate visual similarities as restricted-class evidence.

This is also important ethically. A false positive does not necessarily mean that the model saw nothing meaningful; it may mean that visually similar morphology was given too much weight. Human review would therefore remain essential in any sensitive application.

\section{True negatives: unrestricted dogs correctly classified}

\begin{figure}[p]
    \centering
    \begin{subfigure}[t]{0.84\textwidth}
        \centering
        \includegraphics[width=\linewidth]{TN_curly-coated_retriever_n02099429_1465_gradcam.png}
        \caption{The Curly-Coated Retriever was correctly classified while swimming, despite partial occlusion and an unusual body outline.}
        \label{fig:tn-swimming}
    \end{subfigure}

    \vspace{0.8em}
    \begin{subfigure}[t]{0.84\textwidth}
        \centering
        \includegraphics[width=\linewidth]{TN_greater_swiss_mountain_dog_n02107574_2662_gradcam.png}
        \caption{The Greater Swiss Mountain Dog puppy was correctly classified despite juvenile proportions and an expressive facial pose.}
        \label{fig:tn-puppy}
    \end{subfigure}

    \vspace{0.8em}
    \begin{subfigure}[t]{0.84\textwidth}
        \centering
        \includegraphics[width=\linewidth]{TN_kelpie_n02105412_7810_gradcam.png}
        \caption{The Kelpie was correctly classified, although strong activation falls on the high-contrast pipe and surrounding structure.}
        \label{fig:tn-pipe}
    \end{subfigure}
    \caption{Representative true-negative Grad-CAM explanations.}
    \label{fig:gradcam-tn}
\end{figure}

These examples demonstrate some robustness to unusual conditions. The swimming dog is partly occluded by water, the puppy differs from an adult breed standard, and the Kelpie appears beneath a visually dominant object. Nevertheless, all three were classified correctly.

The pipe example is especially informative. It shows that a correct prediction does not automatically imply a satisfactory explanation. The label is correct, but the model appears to rely partly on an incidental feature of the scene. This distinction between prediction quality and explanation quality is central to the interpretation of post-hoc explanation methods.

\section{True positives: restricted dogs correctly classified}

\begin{figure}[p]
    \centering
    \begin{subfigure}[t]{0.84\textwidth}
        \centering
        \includegraphics[width=\linewidth]{TP_german_shepherd_n02106662_13368_gradcam.png}
        \caption{The German Shepherd was correctly classified as restricted. Activation covers the head and body while extending into the grass.}
        \label{fig:tp-shepherd}
    \end{subfigure}

    \vspace{0.8em}
    \begin{subfigure}[t]{0.84\textwidth}
        \centering
        \includegraphics[width=\linewidth]{TP_rottweiler_n02106550_4962_gradcam.png}
        \caption{The Rottweiler was correctly classified with very high confidence, with substantial overlap across the head, torso, and outline.}
        \label{fig:tp-rottweiler}
    \end{subfigure}

    \vspace{0.8em}
    \begin{subfigure}[t]{0.84\textwidth}
        \centering
        \includegraphics[width=\linewidth]{TP_staffordshire_bullterrier_n02093256_8205_gradcam.png}
        \caption{The Staffordshire Bull Terrier was correctly classified, although the textured sand also contributes to the activation pattern.}
        \label{fig:tp-staffordshire}
    \end{subfigure}
    \caption{Representative true-positive Grad-CAM explanations.}
    \label{fig:gradcam-tp}
\end{figure}

The true positives provide the clearest evidence that the classifier learned useful restricted-class cues. In each example, the heatmap substantially overlaps the dog. The Rottweiler example is particularly persuasive because activation follows the main body region and head. However, none of the maps is perfectly confined to the animal. Grass and sand also receive activation, showing that the final decision may combine morphology with recurring contextual patterns.

\section{Overall interpretation}

Viewed with human eyes, the twelve examples support a nuanced conclusion. It would be inaccurate to claim either that the model fully understands dog breeds or that it merely uses background information. The maps suggest that it often attends to the dog and to plausible breed-related regions, but contextual features also influence some decisions.

The false negatives are associated with difficult presentation: unusual pose or expression, smaller subject scale, diffuse attention, and competing contrast. The false positives often involve unrestricted breeds with visually similar morphology, sometimes combined with environmental texture. The true positives generally show substantial overlap between activation and the dog, while the true negatives demonstrate that the classifier can remain correct under occlusion, juvenile appearance, or unusual setting.

These observations are consistent with the limitations of saliency-based explanation. Heatmaps can be useful for debugging and qualitative inspection, but they should be subjected to cautious interpretation and sanity checking \citep{adebayo2018sanity}. They do not establish that a model reasons in the same way as a human, and they do not make a black-box decision legally or ethically authoritative \citep{rudin2019stop}.

\section{Implications for the project}

The explanation analysis strengthens the project in four ways. First, it confirms that the model frequently attends to the dog rather than only to arbitrary image regions. Second, it reveals contextual dependence that would not be visible from metrics alone. Third, it makes several false-positive and false-negative cases understandable in terms of morphology, pose, scale, expression, occlusion, and contrast. Fourth, it supports the cautious framing of the system as experimental decision support rather than automated legal determination.

A future extension could compare Grad-CAM with alternative explanation methods, quantify localisation against dog bounding boxes, analyse errors by breed, and test whether explanations remain stable under image perturbation. Such work would move the analysis from qualitative inspection towards systematic explanation evaluation.

\section{Chapter conclusion}

Grad-CAM provided a useful qualitative layer of evaluation beyond the confusion matrix. The model often used visually relevant regions on the dog, but it also displayed sensitivity to background texture and high-contrast objects. Correct predictions were not always accompanied by convincing localisation, while some incorrect predictions were understandable because of genuine visual similarity between breeds. The results therefore reinforce the need for human interpretation, transparent limitations, and caution in any sensitive use of the classifier.

\section{Synthesis of the explanation evidence}
Across the twelve cases, the maps provide evidence of partial object localisation rather than a simple background-only shortcut. Several true positives show strong activation over the torso, head and muzzle. Nevertheless, the true-negative Kelpie under the pipe demonstrates that a correct prediction can coexist with substantial attention to contextual structure. Likewise, false negatives involving small scale, unusual expression or patterned backgrounds suggest that the model's representation is sensitive to image composition.

The explanations do not establish a unique set of breed features. Grad-CAM highlights regions, not semantic concepts, and the heatmap resolution is inherited from a late feature tensor. A red area over a dog's head does not reveal whether the network relied on muzzle length, ear position, colour, texture or a correlated combination. More targeted methods could include occlusion tests, counterfactual editing, concept activation vectors and part annotations. Those methods would require additional experimental design and validation.

The explanation study therefore supports three bounded conclusions. First, the model often attends to the animal. Second, attention to the animal does not guarantee a correct decision. Third, context can influence both correct and incorrect predictions. These findings justify human review and robustness testing rather than autonomous use.


\chapter{Discussion, Ethics and Limitations}
\label{chap:discussion}

\section{Answering the research questions}
The first research question asked how effectively frozen ImageNet-pretrained CNNs could classify the constructed binary task. Five of the six architectures exceeded 91\% accuracy, and four exceeded 94\%. Transfer learning was therefore highly effective under the prepared data distribution. The result is consistent with the view that generic visual features can transfer to fine-grained downstream tasks.

The second question concerned architecture selection. InceptionResNetV2 achieved the strongest combination of recall, F1, PR-AUC and MCC. Its advantage was not merely a fraction of accuracy over an otherwise identical model; it produced a notably stronger precision--recall balance than VGG16 and ResNet50. The result is compatible with the architecture's high-capacity multi-scale residual representation, but one experiment cannot prove that architectural design alone caused the difference. Optimisation, input resolution, preprocessing and stochastic variation also contribute.

The third question asked how fine-tuning changed performance. It reduced false positives from 27 to 12 and increased false negatives from 3 to 6. Precision, accuracy, F1 and MCC improved, while recall and the saved AUC values decreased slightly. Fine-tuning therefore shifted the boundary toward fewer positive predictions. This is a meaningful improvement for balanced classification, but not a universal improvement under every possible cost function.

The fourth question concerned explanation. Grad-CAM frequently overlapped the dog and produced plausible regions, especially in clear true-positive cases. It also revealed background sensitivity and did not resolve the semantic basis of errors. The method added diagnostic value but did not convert the network into an intrinsically interpretable model.


\section{Synthesis of the principal finding}
The experiment does not support a simple claim that fine-tuning made every aspect of performance better. The frozen model maximised recall, whereas the fine-tuned model produced the stronger overall balance by sharply reducing false positives. This distinction is important because the preferred model depends on the relative consequences assigned to false negatives and false positives. The final selection therefore reflects a stated metric trade-off, not an assumption that a single headline accuracy fully represents performance.

\section{Why InceptionResNetV2 may have performed strongly}
The task requires evidence at multiple scales. Local colour and texture can distinguish some breeds, but whole-body proportions and head shape require larger receptive fields. Inception branches offer parallel processing paths, while residual shortcuts make very deep feature transformations trainable. ImageNet pretraining also means that the model began with a rich representation of animal parts and textures.

However, architecture narratives can become post-hoc stories. Because InceptionResNetV2 won the benchmark, it is tempting to attribute success confidently to its hybrid design. A stronger causal comparison would control parameter count, input size and training budget or use repeated runs. The present dissertation therefore treats the architectural explanation as plausible, not proven.

\section{Comparison with the proposal and project evolution}
The research proposal envisaged a broader timeline of architectures, systematic hyperparameter search, threshold optimisation and a larger notebook sequence. The completed project narrowed this plan to six established pretrained models, one selected fine-tuning experiment and a fixed Grad-CAM sample. This change improved feasibility and clarity. A broad search with limited compute could have produced many partially controlled comparisons and increased the risk of choosing results after test inspection.

The final design still addresses the proposal's central question: whether architectural progress and transfer learning matter for the restricted/unrestricted task. It also preserves the proposal's concern with interpretability and ethical risk. The divergence lies in scope, not purpose. This is an example of reflective project management: methods were revised in response to resource constraints and the need for a defensible final evaluation.

\section{External validity}
The largest limitation is domain shift. Stanford Dogs images are curated breed examples, often selected because the dog is visible and recognisable. Field photographs may contain multiple dogs, motion blur, poor lighting, occlusion, unusual camera angles and mixed ancestry. A model can perform extremely well within a benchmark and fail when deployed outside it.

Breed labels themselves are not unproblematic ground truth. Visual breed identification by humans is known to be unreliable for mixed dogs, and genetic ancestry does not map perfectly onto appearance. A model trained on labelled purebred examples cannot settle disputes about crossbred animals. The binary label also aggregates distinct breeds into one class and may obscure large variation in per-breed performance.

The unrestricted sample contains twenty-four of the 114 non-positive Stanford categories, not every possible unrestricted breed or cross. The decision boundary is therefore conditional on the selected negative set. Adding visually similar mastiff, bull-type or shepherd breeds could reduce performance. Repeating the experiment with several negative samples would test the stability of the result.

\section{Legal context and category error}
Irish regulations name breeds, strains and crosses, while the model maps a photograph to a dataset-derived binary label. These are not equivalent operations. Legal classification may involve documentary evidence, expert opinion and definitions not captured by a Stanford category. A model score cannot establish that an individual dog falls within a statutory description.

There is also a risk of category error between breed and dangerousness. The project does not predict aggression, bite risk or behaviour. Treating a restricted-class score as a dangerousness score would be scientifically and ethically unjustified. Behaviour depends on many factors outside an image, including socialisation, training, health, environment and immediate context.

\section{Error asymmetry and proportionality}
False positives and false negatives impose different harms. A false negative could fail to identify a dog subject to additional legal controls. A false positive could expose an unrestricted dog and owner to unwarranted intervention. The appropriate balance is not a purely technical choice; it is a policy decision involving proportionality and due process.

The fine-tuned model's twelve false positives and six false negatives may appear small, but deployment scale magnifies error. At ten thousand cases with similar prevalence and performance, errors would number in the hundreds. Moreover, performance could deteriorate under domain shift. High test accuracy should therefore not be used to bypass human judgement or appeal.

A responsible decision-support design would show uncertainty, retain the original image, document the model version, permit expert override and prohibit automatic adverse action. It would also require evaluation across subgroups and image conditions. The present project does not implement such a system; these safeguards define what would be necessary before operational consideration.

\section{Bias and fairness}
Bias can enter through the legal category, dataset and model. The legal list reflects policy choices that are contested in veterinary and animal-welfare debate. The dataset reflects internet availability and breed photography conventions. Some breeds may be photographed in show poses, domestic settings or outdoor working contexts. The model can learn these contextual regularities.

Fairness analysis would require per-breed sensitivity and specificity, confidence distributions, and performance across age, sex, coat variation, image quality and crossbreed status. These attributes are not comprehensively available in Stanford Dogs. The absence of metadata prevents a complete audit. This limitation should not be mistaken for evidence that no disparity exists.

Datasheets and model cards could improve transparency by recording source, selection rules, known exclusions, metrics and intended use \citep{gebru2021datasheets,mitchell2019modelcards}. The reproducibility appendix performs part of this role, but a deployment-grade document would need substantially more governance detail.

\section{Interpretability and contestability}
Grad-CAM provides a visual aid that may help a reviewer identify obvious background dependence. It does not provide a complete explanation of the score and can create false confidence because heatmaps are visually persuasive. Explanation quality must be tested, not assumed \citep{adebayo2018sanity,samek2021explaining}.

In a contested decision, an owner would need more than a coloured overlay. They would need to know the legal basis, model limitations, confidence, alternative evidence and route of appeal. Contestability is an institutional property, not a visualisation feature. This is why post-hoc explanation cannot by itself make a black-box classifier suitable for high-stakes enforcement \citep{rudin2019stop}.

\section{Calibration and threshold selection}
The 0.5 threshold is conventional but not inherently optimal. Calibration asks whether scores correspond to observed frequencies. A model that assigns 0.9 to many cases should be correct roughly 90\% of the time in a calibrated setting. Neural networks can be overconfident, and temperature scaling is one method for post-hoc calibration \citep{guo2017calibration}.

Threshold selection should occur on validation data using a predeclared objective. Possible objectives include a minimum recall requirement, maximum allowable false-positive rate, or explicit cost matrix. The chosen threshold should then be locked before test evaluation. The explanation notebook explored a nearby calibrated threshold, but the main report retains 0.5 for consistency with the established confusion matrix. Future work should separate probability calibration from decision-cost optimisation.

\section{Statistical uncertainty}
The report presents point estimates from one held-out test set. A 98.25\% accuracy is not known without uncertainty. Bootstrap confidence intervals could quantify sampling variation, and McNemar's test could compare paired error patterns between frozen and fine-tuned models. Repeated training seeds would quantify optimisation variance.

These additions were outside the completed pipeline but would materially strengthen the inference. The current result demonstrates that the saved model performed strongly on 1,028 test images. It does not establish a universal performance rate for all dog photographs.

\section{Computational constraints}
Google Colab enabled access to GPU acceleration but imposed runtime and storage constraints. These constraints influenced the choice to benchmark frozen backbones and fine-tune only the selected model. This is methodologically defensible because the project question concerns comparative transfer learning, not exhaustive hyperparameter optimisation.

The large figure files also created Overleaf compilation pressure. The report project therefore uses modular chapter files and supports \texttt{\textbackslash includeonly} for chapter-level builds. The final PDF can be compiled locally from the complete source. This workflow preserves automatic numbering and references while reducing routine edit-compile time.

\section{Strengths of the study}
The study has several strengths. The model zoo uses a common dataset and evaluation pipeline. The test-set class counts are explicit. Model selection is based on several imbalance-aware metrics. The fine-tuning trade-off is reported rather than hidden. Predictions, figures and notebooks are retained. Grad-CAM cases include correct and incorrect outcomes instead of only persuasive successes.

The dissertation also avoids claiming that the legal category is visually natural or behaviourally meaningful. This separation between what the model predicts and what society may infer from it is essential in responsible data analytics.

\section{Future work}
Future research should begin with data, not merely larger models. A stronger dataset would include all relevant breeds and types, crossbred dogs, field images, expert review and documented uncertainty. Exact and perceptual duplicate detection should precede splitting. An external test set from a different source should be reserved for final validation.

Methodologically, future work should report confidence intervals, repeated seeds, calibration, threshold analysis and per-breed metrics. Hard-negative mining could focus training on visually similar unrestricted breeds. Multi-task learning could predict breed before legal mapping, allowing errors to be inspected at the breed level. Open-set recognition or abstention could allow the system to refuse cases unlike its training distribution.

Explanation work should test model sensitivity through occlusion, perturbation and counterfactual examples. A human factors study could examine whether model support improves or harms expert decisions. Most importantly, any operational research would require legal, veterinary, animal-welfare and public consultation rather than treating deployment as a purely engineering question.

\section{Relation to assessment criteria and Level 9 practice}
A capstone artefact is not evaluated solely by whether it operates. The work must demonstrate selection and justification of methodology, integration of theory and practice, independent implementation, coherent reporting, ethical analysis and confident presentation. The project addresses these requirements through the staged notebook pipeline, architecture comparison, error analysis and explicit limitations.

Critical reflection is particularly important. Several decisions changed during the project: the focus moved away from unavailable XL Bully data; the architecture set was narrowed; exhaustive hyperparameter search was not pursued; and the explanation method was stabilised through multiple notebook revisions. Reporting these changes demonstrates management of research constraints. Concealing them would create a falsely linear account of development.

\section{Alternative modelling formulations}
The direct binary classifier is not the only possible formulation. A multiclass model could first predict one of the thirty selected breeds and then map the predicted breed to legal status. This would provide breed-level transparency and reveal which restricted breeds drive performance. However, multiclass training would require distinguishing many more boundaries and could reduce binary accuracy.

A hierarchical model could first distinguish broad morphological groups and then make a finer prediction. Metric learning could embed visually similar dogs near one another and support nearest-neighbour inspection. An ensemble could average the strongest architectures. Each approach introduces complexity and new validation requirements. The direct binary classifier was appropriate for a bounded comparative study, but it should not be mistaken for the only defensible design.

\section{Why an ensemble was not selected}
InceptionV3, Xception, InceptionResNetV2 and NASNetMobile all performed strongly, so an ensemble might reduce uncorrelated errors. Probability averaging could improve ranking and robustness. The likely gain must be weighed against additional inference cost, storage and interpretability complexity.

Because the selected fine-tuned single model already made eighteen test errors, an ensemble could at most improve a relatively small absolute number of cases on this split. Without an independent validation design, ensemble weights could also overfit. The project therefore favoured a single well-documented model. Future work could test whether the models' errors are sufficiently diverse to justify ensembling.

\section{Human-in-the-loop design}
A responsible support tool would not present a binary verdict in isolation. It could display the score, uncertainty interval, nearest training examples, model limitations and an abstention warning when the image is out of distribution. A trained reviewer would combine this information with documentation, physical examination and other evidence.

Human involvement does not automatically remove bias. Reviewers can over-trust algorithmic output, a phenomenon often described as automation bias. Interface design and training would need to emphasise that the score is advisory. Overrides and disagreements should be recorded so that systematic model failures can be audited.

\section{Data protection and image governance}
The Stanford dataset is publicly available for research, and the project does not process owner identities. Nevertheless, an operational system could receive photographs containing people, addresses, vehicle registrations or location metadata. Such a system would process personal data and would require a lawful basis, purpose limitation, minimisation, retention controls and security safeguards.

The research artefact does not implement these operational functions. The ethical analysis includes them because deployment changes the data-protection profile. Moving a model from a benchmark notebook into a public-service workflow is not a neutral technical step; it creates new governance obligations.

\section{Animal welfare perspective}
Errors affect animals as well as owners. False identification can contribute to restrictive handling, prolonged kennelling or adverse legal outcomes. Even correct breed identification does not establish that a dog poses a behavioural threat. A welfare-aware system should therefore avoid language that equates restricted status with individual dangerousness.

The model could also have beneficial uses in a carefully designed setting, such as prompting a second expert review rather than triggering action. Whether benefits exceed harms depends on the surrounding institution. The dissertation cannot answer that policy question from image metrics alone.

\section{Reflexive account of the research process}
The project involved repeated negotiation between ambition and reproducibility. It was tempting to add more models, searches and diagrams, but every addition consumes time and creates another opportunity for uncontrolled comparison. The final pipeline became stronger when it was simplified around explicit research questions.

The work also improved through attention to presentation. Confusion matrices were standardised, tables were reformatted, explanation cases were arranged vertically for inspection, and chapter source was modularised to overcome compilation constraints. These are not merely cosmetic changes. Clear presentation supports scrutiny and reduces the chance that an important trade-off is hidden by layout.

The most important learning outcome was that a high score is the beginning of interpretation rather than the end. The final model's 98.25\% accuracy initially appears decisive. Examining the eighteen errors, the threshold, the missing breeds and the Grad-CAM maps produces a much more qualified understanding. That movement from headline metric to bounded claim is central to advanced data-analytics practice.


\chapter{Conclusion}
\label{chap:conclusion}

This dissertation benchmarked and critically evaluated a reproducible ImageNet transfer-learning pipeline for binary restricted versus unrestricted dog-image classification. Six ImageNet-pretrained CNN architectures were compared as frozen feature extractors using a common Stanford Dogs-derived dataset. InceptionResNetV2 achieved the strongest frozen performance and was selected for limited fine-tuning.

The frozen model achieved 97.08\% accuracy and 98.42\% recall, with 27 false positives and 3 false negatives. Fine-tuning improved accuracy to 98.25\%, precision to 93.88\%, F1 to 95.34\% and MCC to 0.9428. It reduced false positives to 12 while increasing false negatives to 6. The central empirical conclusion is therefore nuanced: fine-tuning improved overall balance and positive predictive reliability, but the frozen model remained more sensitive to restricted examples.

The architecture comparison supports the usefulness of modern multi-scale and residual representations for this controlled task. It does not establish that architecture alone caused the observed ranking, nor that the model will retain its performance outside the prepared split. The dataset is curated, incomplete with respect to the legal list, and poorly suited to crossbred or ambiguous real-world cases.

Grad-CAM added a qualitative diagnostic layer. The selected maps often overlapped the animal, but correct predictions sometimes attended to context and incorrect predictions could still appear visually plausible. The explanations therefore support inspection without proving human-like reasoning or legal validity.

The project meets its aim as a controlled academic benchmark and proof of concept. It demonstrates dataset engineering, transfer-learning comparison, fine-tuning, metric-aware evaluation, checkpointing, explanation and critical reflection. It does not demonstrate an autonomous legal classifier. Photograph-only output should never be treated as a substitute for expert evidence, due process or behavioural assessment.

The most valuable next step is not simply to increase model capacity. It is to improve external validity through representative field data, crossbreed coverage, duplicate control, calibration, uncertainty estimation, subgroup analysis and independent testing. Any movement toward operational use would require human-in-the-loop design, transparent documentation, appeal and governance. The technical result is strong; the responsible conclusion remains deliberately bounded.

\section{Final reflection}
The project also demonstrates why evaluation quality matters as much as model construction. A less critical account could have reported only the 98.25\% accuracy and described the experiment as solved. The fuller analysis shows that the same model made twelve false-positive and six false-negative decisions, that precision depends on prevalence, that the positive class is legally rather than visually defined, and that the source dataset omits important breeds and mixed dogs. These qualifications do not diminish the technical achievement. They identify exactly what has and has not been demonstrated.

The comparison of six backbones provided a defensible basis for model selection, while the frozen-versus-fine-tuned comparison showed that additional training changes the operating trade-off. The explanation analysis further showed that correct classification and satisfactory visual reasoning are not equivalent. Together, these findings support a broader lesson: responsible machine learning requires an evidence chain from data construction through model training, evaluation, explanation and governance.

For the MSc capstone, the completed artefact is therefore both the trained classifier and the documented research process around it. The notebooks, saved predictions, figures and chapter source make the decision path inspectable. The report's final position is intentionally balanced. CNN transfer learning can classify the prepared images with high accuracy and can provide useful comparative insight into architecture and fine-tuning. It cannot, on the evidence presented here, determine the legal identity, ancestry or dangerousness of an arbitrary dog. Maintaining that boundary is central to the validity of the work.


% ============================================================
% APPENDICES
% ============================================================

\appendix

\chapter{Reproducibility Artefacts}
\label{app:reproducibility}

The project archive contains the four principal notebooks, saved figures, result tables and LaTeX source. The notebooks are independent of a previous Google Colab runtime but retain the dependencies of the experimental pipeline: data preparation precedes frozen-model benchmarking; the selected model checkpoint is required for fine-tuning; and the model explainability notebook requires the trained model and test dataset.

\begin{itemize}
    \item \texttt{01\_Data\_Preparation.ipynb}: dataset acquisition, class mapping, validation and split construction;
    
    \item \texttt{02\_Frozen\_Model\_Benchmarking.ipynb}: frozen transfer-learning benchmark across the six candidate CNN architectures;
    
    \item \texttt{03\_Selected\_Model\_Fine\_Tuning.ipynb}: fine-tuning and evaluation of the selected InceptionResNetV2 model;
    
    \item \texttt{04\_Model\_Explainability.ipynb}: prediction verification, error analysis and Grad-CAM visualisation.
\end{itemize}

The report is divided into separate chapter files and assembled through \texttt{report.tex}. GitHub Actions compiles the report, poster, presentation and slides whenever their source or supporting figures change. A complete build is performed so that the table of contents, cross-references, bibliography and chapter auxiliary files remain current without depending on Overleaf.


\chapter{Detailed Architecture Summary}
\label{app:architecture}
\begin{longtable}{p{2.6cm}p{3.2cm}p{4.0cm}p{3.2cm}}
\caption{Block-level summary of the evaluated architectures.}\label{tab:blocksummary}\\
\toprule
Model & Stem and reduction & Principal blocks and kernel patterns & Final representation \\
\midrule
\endfirsthead
\toprule
Model & Stem and reduction & Principal blocks and kernel patterns & Final representation \\
\midrule
\endhead
VGG16 & Repeated convolution followed by $2\times2$ max pooling & Two or three $3\times3$ convolutions per block; channel depth increases from 64 to 512 & $7\times7\times512$ before the original dense top; global average pooling used in this project \\
ResNet50 & $7\times7$ convolution and max pooling; strided residual stages & Bottleneck blocks use $1\times1$, $3\times3$, $1\times1$ convolutions with identity or projection shortcuts & 2048-channel terminal feature tensor followed by global average pooling \\
InceptionV3 & Factorised convolutional stem and reduction modules & Parallel branches using $1\times1$, $3\times3$, and asymmetric $1\times7$/$7\times1$ factorisations & 2048-channel representation before pooling \\
Xception & Two-convolution entry stem and strided entry flow & Depthwise $3\times3$ spatial convolution plus pointwise $1\times1$ channel mixing, arranged in residual blocks & 2048-channel exit-flow representation \\
InceptionResNetV2 & Multi-convolution stem, Reduction-A and Reduction-B & Repeated Inception-ResNet A, B and C blocks with parallel kernels, projection and scaled residual addition & 1536-channel terminal representation before pooling \\
NASNetMobile & Convolutional stem and learned reduction cells & Repeated normal cells preserve resolution; reduction cells downsample; separable convolutions selected by architecture search & Approximately 1056 channels before global pooling \\
\bottomrule
\end{longtable}


% ============================================================
% REFERENCES
% ============================================================

\clearpage

\addcontentsline{toc}{chapter}{References}

\bibliographystyle{apalike}

\bibliography{references}

% ============================================================
% END DOCUMENT
% ============================================================

\end{document}